\documentclass[12pt]{article}

\usepackage{cmap}
\usepackage[T2A]{fontenc}
\usepackage[utf8]{inputenc}
\usepackage[english,russian]{babel}
\usepackage{tabularx}
\usepackage{graphicx}
\usepackage{amsmath}
\usepackage{indentfirst}
\usepackage{tikz}
\graphicspath{ {./images/} }

\newcommand{\ornament}[1]{%
  \begin{tikzpicture}[baseline]
    \draw[line width=0.4pt] (0,0) -- (#1,0);
    \node at (#1+0.45,0) {\footnotesize$\diamond$};
    \draw[line width=0.4pt] (#1+0.9,0) -- (2*#1+0.9,0);
  \end{tikzpicture}%
}

\usepackage{xcolor}
\definecolor{pastelblue}{rgb}{0.0, 0.2, 0.4}
\usepackage[colorlinks=true, linkcolor=pastelblue]{hyperref}

\title{Конспект лекций Михаила Густокашина по дисциплине «Алгоритмы и структуры данных 2»}
\author{От Добросовестной Слушательницы Ли Людмилы БПМИИ247}
\date{2025 год}

\begin{document}
\maketitle

\begin{flushright}
<<Ваше будущее --- это перекладывать джейсоны>> \\
--- Михаил Сергеевич Густокашин \\ [3ex]
<<Вы не могли задать такой глупый вопрос>> \\
--- Михаил Николаевич Вялый \\ [3ex]
<<Если вас отчислят по алгосам, хоть в художку пойдёте>> \\
--- Борис Васильевич Галицкий

\end{flushright}

\begin{figure}[h]
\centering

\includegraphics[width=0.5\linewidth]{gust1.jpg}
\end{figure}


\newpage

\tableofcontents


\newpage


\section*{Лекция 1: Префикс-функция}
\addcontentsline{toc}{section}{Лекция 1: Префикс-функция} 

Поговорим об алгоритмах на строках. Мы рассмотрим такие алгоритмы, которые в основном помогают решать задачи по типу поиска подстроки в строке. Это фундаментальные вещи. 

\subsection*{Префикс-функция}
\addcontentsline{toc}{subsection}{Префикс-функция}

\textit{Префикс функция определена для каждой позиции в строке и обозначает максимальную длину суффикса подстроки, которые заканчиваются в одном месте, совпадающем с префиксом}. Пусть дана строка aababaаaba и нам надо посчитать значение префикс функции для каждой позиции. Какая максимальная по длине подстрока заканчивающаяся на первом элементе, но не совпадающая со всем началом строки, совпадает с началом? Никакая: там один символ. Подходит только пустая строка. Дальше строка аа. У нее есть 2 суффикса: а и аа. Тот суффикс, который аа, он совпадает со всей строкой, поэтому он нам не интересен. Мы рассматриваем только те суффиксы, которые не совпадают со всей строкой. В нашем случае это буковка а. Мы должны понять, какой из наших суффиксов (в данном случае он единственный) максимально совпадает по длине с префиксом всей строки: буква а совпадает с префиксом из одной буквы а. Длина 1, его и записываем. Теперь смотрим на b. Какие здесь есть суффиксы? b и ba. Соответствующие по длине префиксы --- это а и аа. Какой максимальный по длине суффикс совпадает с префиксом? Никакой, пишем 0. Дальше: a, ba и aba. Какой максимально по длине совпадает? a, значит пишем 1.

\begin{table}[htbp]
  \centering
  \resizebox{\textwidth}{!}{
    \label{tab:example}
    \begin{tabular}{|c|c|c|c|c|c|c|c|c|c|}
      \hline
      a & aa & aab & aaba & aabab & aababa & aababaа & aababaаa & aababaаab & aababaаaba \\
      \hline
      0 & 1 & 0 & 1 & 0 & 1 & 2 & 2 & 3 & 4\\
      \hline
    \end{tabular}
  }  
\end{table}

Допустим, мы хотим найти все вхождения строки аа (Aлгоритм Кнута-Морриса-Пратта). Мы пишем аа, затем символ-разделитель (\#), потом строку, в которой мы ищем нашу подстроку, а дальше считаем префикс функцию.

\begin{figure}[h]
\centering
\includegraphics[width=0.4\linewidth]{aaapic1.png}
\end{figure}

Теперь научимся это считать. У наивного алгоритма сложность будет кубическая, а бинпоиском здесь воспользоваться не получится.

Как у нас ведет себя префикс-функция? Заметим, что в какой-то момент она начала практически возрастать. А на сколько максимум префикс-функции может сдвинуться при сдвиге при добавлении одного символа? На 1. А уменьшиться она может на сколько угодно. 

\begin{figure}[h]
\centering
\includegraphics[width=0.55\linewidth]{aaapic2.png}
\end{figure}

Теперь рассмотрим пример поинтереснее. У нас уже есть последний символ, совпадающий с началом строки, то бишь четвертая буковка а, у которой значение равно 1. Может ли значение префикс-функции последней буковки а стать 5? Нет, ведь стоит буква с. А 4? Оно могло бы стать равной 4 только в ситуации, если бы у четвёртой буквы было бы значение функции 3. Потому что, чтобы значение префикс функции стало равно 4, нужно чтобы aba совпадало с тремя символами префикса. Но мы можем быстро понять, совпадают ли эти символы с префиксами. А как это быстро сделать? Посмотреть на значение функции под четвертой буковкой а: если бы оно было равно 4, то это значит, что последние три символа из начала совпадают с тремя первыми символами. В нашем примере оно не равно 3, значит, aba не совпадает с первыми тремя символами, а значит, значение префикс-функции последнего элемента не может быть 4. А может быть, оно равно 3? Смотрим: ba должно совпадать с первыми двумя символами. Нет! Может быть, оно равно двум? Нужно, чтобы предпоследняя буковка а совпадала с началом строки. Совпадает! Значит, у подстроки, заканчивающейся на последнюю букву а, значение префикс-функции будет 2.

\begin{figure}[h]
\centering
\includegraphics[width=0.4\linewidth]{aaapic3.png}
\end{figure}

Но мы сделали очень много лишних действий. Если у нас есть число, с которым сравниваем (это значение под четвертым символом строки), то берем его, а не перебираем все числа подряд.

Рассмотрим еще один пример. Первое значение префикс-функции всегда 0. Как мы для следующей позиции определяем значение префикс-функции? Смотрим, совпадает ли этот символ с символом, который стоит после префикса, длина которого равна значению префикс-функции в предыдущей позиции. \textbf{Нумерация у нас с единицы.} Далее мы знаем, что предыдущая префикс-функция была 1. Теперь пытаемся приписать элемент к максимально совпадающему суффиксу после этого символа. Какой максимальный суффикс, не считая а, совпадает с префиксом всей строки? Да никакой, длиной 0. И мы пытаемся к строке длины 0 приписать букву b. Успешно сравнили? Нет, значит ноль. Встретили букву а. Пытаемся приписать ее в предыдущей строке. Приписываем к строке длинной 0 и сравниваем со следующим за префиксом длины 0. Буквы а совпали. Пишем число на один больше. Дальше буква b. До этого стояла единица, смотрим что до неё, приписываем b к нулю и сравниваем: безуспешно. То же самое со следующей буквой а. Следующую букву а мы приписываем к а, получаем значение 2. Следующая буква а так же будет иметь значение два, потому что мы пытались приписать третью букву а --- не получилось, но мы смотрим, а какой максимальный суффикс строки ааа совпадает с началом строки.

На 14-ом месте у нас стоит буква b, но мы не сразу пишем ноль. Перед b стоит a, значение которого 7. Значит, мы переходим к седьмому элементу и смотрим какое значение у него --- в нашем случае это 2. Теперь возвращаемся в начало строки, а именно к элементу под номером 2 и приписываем к нему b. Совпадает! Значит буква b под 14-ым номером имеет значение 3.

\begin{figure}[h]
\centering
\includegraphics[width=0.55\linewidth]{aaapic4.png}
\end{figure}

Теперь напишем код:

\begin{verbatim}
max pref = p[i - 1]
while max pref > 0 && s[i] != s[max pref + 1]:
    max pref = p[max pref]
p[i] = max pref + 1
\end{verbatim}

Не забудьте учесть случай, когда функция обращается в ноль! Почему он работает за линейное время? У нас есть переменная max pref: она может либо увеличиться на единичку, либо уменьшиться. Всего $N$ шагов. Сколько действий может произойти с этой переменной? $N$, поскольку количество уменьшений не может превосходить количество увеличений, а кол-во увелечений максимально $N$. Поэтому префикс-функция считается за линейное время. 

Представим, что у вас появилась потребность написать онлайн-алгоритм, то есть которая работает в реальном времени. В таком случае мы не можем хранить всю эту последовательность. Если мы говорим об алгоритме Кнута-Морриса-Пратта, то мы храним только первые элементы. У нас есть какая подстрока, разделитель, а потом идёт длинная строка. На самом деле нам нужно только те элементы, которые находятся перед разделителем, потому что он \textit{уникальный}. Нам нужно хранить только первый кусочек и последний элемент. 

А как сделать этот алгоритм за О(1) в целом? Мы хотим для каждого значения префикс-функции и для каждой возможной последующей буквы сделать такую табличку:

\begin{center}
\begin{tabular}{c c c c c}
 - & 0 & 1 & 2 & 3\\
 a & 1 & 1 & 3 & 1 \\
 b & ... & & \\
 c & \\
\end{tabular}
\end{center}

И таким образом вы уже знаете буквы и значения префикс-функции, и уже попадаете за О(1). 

\subsection*{Z-функция}
\addcontentsline{toc}{subsection}{Z-функция}

Нам нужно для каждой позиции определить максимальную длину подстроки, начинающуюся с этой позиции и совпадающей с префиксом всей строки. Если у префикс-функции \textit{заканчивалась} строка в этой позиции, то у нас какая строка \textit{начиная} с этого места совпадает с префиксом.

\begin{figure}[h]
\centering
\includegraphics[width=0.6\linewidth]{aaapic6.png}
\end{figure}

Сейчас мы будем это учиться считать за линейное время. Мы будем хранить так называемый $Z$-блок --- это максимальная подстрока, совпадающая с префиксом, то есть мы поддерживаем $Z$-блок правым концом (см. картинку снизу). Например, для позиции 7 --- он будет длиной 7. В каждый момент времени мы поддерживаем отрезок с самым правым концом. Что мы узнаем благодаря этому? Что этот кусок совпадает с префиксом всей строки, и это сейчас позволит нам более быстро пересчитывать $Z$-функцию. 

Мы хотим посчитать для следующей позиции --- восьмёрки. Она у нас покрывается $Z$-блоком. Что мы можем сделать? Мы можем посмотреть на второй символ! Мы знаем, что эти семь символов совпадают с первыми семью цифрами всей строки. Смотрим на второй элемент и переписываем его (см. картинку снизу). Это работает до тех пор, пока не окажется, что значение $Z$-функции вылезает за пределы нашего отрезка. И как только мы натыкаемся на блок неизвестности (то есть на конец $Z$-блока), мы начинаем по-честному считать.

\begin{figure}[h]
\centering
\includegraphics[width=0.6\linewidth]{aaapic7.png}
\end{figure}

Почему это очень быстро считается? Если мы оказались в блоке, то мы просто берем значение за О(1). Если мы превысили правую границу $Z$-блока, то мы будем идти, начиная с границы блока до первого несовпадения, и наш $Z$-блок новый будет расширяться туда, либо не будет, если мы попали на несовпадение. То есть на каждом шаге мы понимаем, что либо мы сейчас $Z$-блок не расширим, либо расширяется блок на каждый символ, тратя О(1). Но в таком случае мы не можем написать эффективный онлайн-алгоритм. 



\newpage
\section*{Лекция 2: Ахо-Корасик}
\addcontentsline{toc}{section}{Лекция 2: Ахо-Корасик} 

Сегодня мы будем заниматься примерно тем же, что и на прошлом занятии.

\subsection*{Бор (Trie)}
\addcontentsline{toc}{subsection}{Бор (Trie)}

У нас такая задача: предположим у нас есть какой-то словарь слов и нам нужно как-то реагировать и слово быстро проверять. Наш словарь состоит из следующих слов: a, abac, acb, caba, и мы хотим проверять слово. Строим наш бор: у нас есть корень дерева, в котором нет ничего, и мы начинаем потихонечку строить этот бор. Вот из пустой строки у нас есть возможный переход по букве а. Когда слово заканчивается, мы ставим \textit{крестик} ($\times$). Это называется признак \textit{терминальности}. Теперь создаём бор для слова abac. Создаём стрелку для b, a, c. Теперь создаём бор для слова acb. От a создаём ветку c, b. Аналогично с caba. (см. картинку)

\begin{figure}[h]
\centering
\includegraphics[width=0.7\linewidth]{aaapic8.png}
\end{figure}

Как это писать? В идеале нам нужно, чтобы переход делался за О(1). То есть нам понадобится либо какой-то массив с переходами (если это латинский алфавит, можно порядковый номер буквы считать), либо мы можем использовать словарь (map, unordered\_map возможно, с переходами: ключ --- это char, значение --- прямо нода, не указатель). Unordered\_map содержит в себе словарик с возможными переходами из неё и булевую переменную (признак терминальности). 

Представим, что нам нужно динамически добавлять и удалять слова из нашего словаря. Добавлять-то мы поняли как, а как удалять? Для каждой вершины нам нужно завести счётчик, чтобы знать, сколько слов через неё проходит. Если счётчик стал нулём после вычитания единички (потому что мы удалили слово, которое проходило через эту вершину), то вершину можно смело удалять. 

Мы введем новое обозначение: trie(v, c), где v --- это текущая вершина, с --- символ. Соответственно, данная функция показывает, в какую вершину мы попадаем, переходя из вершины v по символу c. Это всё прекрасно до тех пор, пока у вас в тексте есть пробелы. Этого может и не быть.

\subsection*{Ахо --- Корасик}
\addcontentsline{toc}{subsection}{Ахо --- Корасик}

Представим, что у нас колбаса из символов и нам нужно реагировать каждый раз, когда мы видим в ней определенную строку. Здесь бор не поможет. Для этого в 1975 году придумали алгоритм под названием Ахо-Корасик (это фамилии, склонять первую часть нельзя!). Берём тот же бор и строим на нём \textit{суффиксную ссылку}. А именно для каждой вершины мы должны научиться переходить к самой длинной подстроке, которая у нас является суффиксом нашей текущей строки. Давайте договоримся, что рисовать те суффиксные ссылки, которые ведут в корень дерева, я не буду. То есть ни одного такого суффикса нет, и поэтому если зеленая стрелочка не ведет из вершины никуда, значит, она идёт в корень нашего бора. 

Например, строка ab. У нее есть какой-то суффикс, который есть в боре? Ну, единственный суффикс здесь b, если тут такая строка? Нет, значит, я не рисую ничего. Далее aba. Если ли какой-нибудь суффикс здесь, который есть в боре? Да, а, поэтому рисуем зеленую стрелочку, направленную на конец стрелочки а. Напоминаю, что те вершины, из которых ничего не идёт, ведут в корень, а также о том, что мы ищем самые длинные суффиксы. В abac это ac, поэтому проводим зеленую стрелочку в c. И так далее. (см. картинку снизу)

И зная суффиксные ссылки, мы можем делать примерно то же, что мы и делали с префикс-функциями. Допустим, у нас встретилась строка cabac. Мы идём с корня, то есть самого начала строки. Далее мы обрабатываем caba, фиксируем, а дальше переходим по суффиксной ссылке. Переходим по ней и оказываемся в конце. Радуемся! Когда мы не можем пройти по букве дальше, мы идём по суффиксной ссылке. А если бы у нас было не cabac, а cabab? Мы идём по caba, зафиксировали, дальше нам надо идти по букве b, но по букве b нет перехода. Идём по суффиксной ссылке. Опять нет перехода. Идём снова по суффиксной ссылке и обнаруживаем букву b. Мы переходим по букве b и оказываемся в соответствующем состоянии. А если там было бы не cabab, а cababac, то мы бы прошлись еще и по a и c. Зачем нам нужно поддерживать состояния? Состояния --- это максимальный префикс, какой-либо из строк, который у нас сейчас накоплен в некотором буфере. Это нужно, чтобы каждый раз заново не ходить от корня. 

\begin{figure}[h]
\centering
\includegraphics[width=1\linewidth]{aaapic9.png}
\end{figure}

Почему префикс-функция работала быстро? Потому что мы по каждой встретившейся буковке либо движемся вправо по нашему бору, либо идём к уменьшению, ведь по суффиксной ссылке мы всегда к более короткой строке проходим. И работает это всё за линейное время. 

У нас есть три проблемы: мы как-то плохо строили этот автомат, мы строим его медленно и работает это всё амортизированно за О(1). 

На самом деле все эти проблемы решаются одним способом, но начнём мы с производительности, то есть работы автомата. Мы хотим сделать автомат Ахо-Корасик, который будет обозначать следующее. Я обозначу это как go(v, c), где v --- вершина, c --- символ. То есть куда мы попадаем из вершины v по символу c. Если такой переход у нас в боре есть, то всё прекрасно. А если нет c в боре? Мы хотим наш автомат за О(1) сделать. Куда мы попадём из вершины a по букве c? В тупик. Переход у нас есть, поэтому мы рисуем красным маркером стрелочку в конец. Если нам нужно пройти по b, то красную стрелочку рисуем к b. Если нам нужно пройти по a, то мы сначала два раза проходим по суффиксной ссылке, а далее переходим в корень, потому что если стрелочки нет, то она ведет в корень. 

\begin{figure}[h]
\centering
\includegraphics[width=1\linewidth]{aaapic10.png}
\end{figure}

Чем trie(v, c) отличается от go(v, c)? Если мы не могли в боре куда-то пройти, то мы расстраивались, а go никогда не будет расстраивать, он будет куда-то идти. То есть go --- это автомат, который по состоянию и очередному символу переходит в какое-то другое состояние.

\begin{verbatim}
go (v, c) --- это:
trie(v, c), если с есть в вер. v бора
иначе go(suf(v), c)
\end{verbatim}

Это будет работать, но вообще это лучше решать обходом в ширину от корня нашего дерева, то есть обрабатывать такой волной. 

Мы решили только одну из трёх проблем. Давайте теперь научимся считать суффиксные ссылки быстро. 

\begin{verbatim}
suf (v) = go(suf(parent_v), parent_c)
\end{verbatim}

Как это работает? Допустим, мы находимся в строке abac и хотим посчитать куда ведёт суффиксная ссылка для элемента c. Мы берем aba и находим куда ведет его суффиксная ссылка. Допустим, у нас есть ba. Мы попадаем в строку ba и делаем шажок по букве c. И если у вас в боре была строка bac, то суффиксная ссылка будет вести в строку bac

\begin{figure}[h]
\centering
\includegraphics[width=0.65\linewidth]{aaapic11.png}
\caption{Как работает нахождение суффиксной ссылки через go}
\end{figure}

\begin{figure}[h]
\centering
\includegraphics[width=0.65\linewidth]{aaapic12.png}
\caption{Как работает нахождение суффиксной ссылки через trie (много шагов)}
\end{figure}

Мы научились строить неработающий автомат. Почему неработающий? Он как-то работает, но не срабатывает на этой букве а несчастной. Берем синий маркер! И сделаем \textit{терминальную ссылку}. Она будет вести во все терминалы, которые образовались в данный момент. То есть мы пришли в какое-то состояние и у нас синяя стрелочка, то есть терминальная ссылка, ведет в то слово, которое сейчас закончилось. Как нам определить, какой самый длинный суффикс является терминалом? Очень просто: давайте пропрыгаем по всем суффиксным ссылкам, и если у нас встретится терминал (терминальный символ), то мы зафиксируем это. Предположим, мы пришла в ca от корня. Каждый раз когда мы приходим куда-то, мы прыгаем по всем суффиксным ссылкам, пока не дойдём до корня, и если у нас где-то встретился терминал, то фиксируем его вхождение.

Почему это работает? Потому что суффиксная ссылка ведет на самый длинный суффикс, который встречается в нашем боре. И если этот суффикс, соответственно пропрыгав по всем терминальным ссылкам до корня, мы переберем все суффиксы, которые встречаются в нашем боре. И если среди суффиксов окажется терминал, мы должны его зафиксировать. Но это занимает много шагов. Давайте сократим их: зачем нам прыгать по нетерминальным вершинам? Давайте из каждой вершины, вместо того чтобы на самый большой суффикс вести, сразу будем вести терминальную ссылку на самый большой терминальный суффикс и из него без промежуточных сразу в большой терминальный. Бегаем по всем суффиксным ссылкам, в надежде обнаружить терминал. 

\begin{figure}[h]
\centering
\includegraphics[width=0.75\linewidth]{aaapic13.png}
\caption{Бегаем по всем суффиксным ссылкам, в надежде обнаружить терминал}
\end{figure}

Как его эффективно написать?

\begin{verbatim}
term(v) --- это:
suf(v), если suf(v) - терминал
иначе term(suf(v))
\end{verbatim}

По терминальной ссылке определяем вхождение символов и суффиксов.

\begin{figure}[h]
\centering
\includegraphics[width=1\linewidth]{aaapic14.jpeg}
\caption{Мем с доски}
\end{figure}

\newpage
\section*{Лекция 3: Алгоритмы сжатия и кодирования}
\addcontentsline{toc}{section}{Лекция 3: Алгоритмы сжатия и кодирования} 

Сегодня мы будем учиться сжимать данные. Рассмотрим, как работает zip.

\subsection*{Как работает zip? Алгоритм LZW}
\addcontentsline{toc}{subsection}{Как работает zip? Алгоритм LZW}
Первое, о чём мы начнём говорить, --- это алгоритм LZW. На самом деле это целое семейство алгоритмов. Рассмотрим на примере. Возьмём строку ababacabaca. Не забываем, что нам нужно не только запаковать, но и распаковать потом. Сначала у нас в словаре содержатся все байты возможные, то есть у нас есть нулевой байт, так далее, затем код буковки a, так далее и, соответственно, до 255. Договоримся так, хотя это очень зависит от реализации. Это начальное состояние. Потом мы берем очередную букву из строки и проверяем, есть ли такое слово в словаре. Мы взяли первую букву a, и она оказалась в нашем словаре. Берем следующую букву. Есть слово ab в словарике? Нет, значит, мы его туда добавляем. Каждое следующее слово состоит из предыдущего слова и одной дописанной буквы. И так далее.

\begin{figure}[h]
\centering
\includegraphics[width=0.5\linewidth]{aaapic15.png}
\caption{Словарь}
\end{figure}

Как это реализовывать? Вы можете подстроки использовать как ключи вашего словаря (unordered\_map, если вы не гоняетесь за скоростью). 

Подытожим: берем очередную буковку, если такое есть в словаре, берем следующую буковку и так далее, пока не окажется, что его нет словаре. Тогда выписываем номер последнего слова, которое всё-таки было и очередную добавленную букву. 

Теперь учимся распаковывать! На руках у нас есть следующая строка:

\[
'a'b256a'c'a'b'a258b
\]

Мы договорились, что у нас есть словарь, поэтому воспользуемся им. Зная, что первые две буквы это a и b, поэтому мы знаем, что в момент создания 256, наш алгоритм упаковки создавал новое слово в словаре. Из чего он состоит? Из предыдущего слова и одной новое буквы. Мы их и пишем. Записываем ababacaba258b. Параллельно записывая каждую букву, составляем словарь.

Этот алгоритм хорош тем, что вместе с ним не нужно передавать словарь, а нужно передавать только упакованную строку, а словарь восстанавливается по ходу распаковки.

\begin{figure}[h]
\centering
\includegraphics[width=0.62\linewidth]{aaapic16.png}
\end{figure}

Уберем букву b в конце. Как это упаковать тогда? Да, действительно пришлось бы повторно использовать слово. Вот у нас уже есть слово ca, мы идём-идём, пришли к тому, что нужна буква, но буквы нет. Поэтому мы делаем копию слова и заново создаём слово ca. 

\begin{figure}[h]
\centering
\includegraphics[width=0.62\linewidth]{aaapic17.png}
\end{figure}

Сколько бит мы тратим на запись номера слова? на 256 мы тратим 9 бит, на буковку а в начале --- 8, а как только наш словарь достигнет размера 511, начиная с 512 места мы будем тратить 10 бит на запись номера слова. Затем соответственно 11, 12 и т.д. А какой должен быть размер у словаря? Нам нужно придать ему адаптивности. Наш словарь будет выглядеть в виде дерева: у него есть некий корень, из него торчат разные буковки (255 возможных байтов в нашем способе кодирования). Нам интересны буковки a, b и c (см. рисунок снизу)

\begin{figure}[h]
\centering
\includegraphics[width=1\linewidth]{aaapic18.png}
\caption{Запись строки в виде бора}
\end{figure}

Теперь у нас может сложиться такая ситуация, что место в словаре кончилось и нужно кого-нибудь удалить. Можем ли мы удалить слово ab? Нет, потому что на него ссылается слово aba. В этом дереве можно удалять только листья. Но на самом деле у нас есть еще куча листьев, которые соответственно для неиспользованных символов, но их трогать нельзя. Для каждого листа мы храним момент последнего использования, и вот самый старый из них мы можем удалить, если наш словарь переполнился. Можно также хранить сколько раз мы проходили через этот лист. Но тут возникает проблема: каждый лист мы посещаем по одному разу. Но такой метод имеет место быть, если вершина стала листом после удаления какой-то ветки и тогда получается, что она посещалась больше одного раза изначально.

\subsection*{Алгоритм Хаффмана}
\addcontentsline{toc}{subsection}{Алгоритм Хаффмана}

Предположим, что в нашем тексте 60 буковок а, 40 буковок b, 20 c и 10 d. Короткий код присваиваем букве а, а самый длинный --- букве d. В чём заключается идея алгоритма: мы берем две наиболее редкие буковки (это c и d в нашем случае) и объединяем их в один узел с суммарным весом. Суммарный их вес --- 30. Далее минимальные веса имеют значения 30 и 40. Объединяем их и получаем 70. Объединяем в последний раз: получаем суммарный вес 130. Это нужно, чтобы маловстречающимся буквам сопоставить длинный код, а частовстречающимся --- короткий. Справа пишем 0, а слева 1. (см. картинку снизу)

\begin{figure}[h]
\centering
\includegraphics[width=0.7\linewidth]{aaapic19.png}
\caption{Сопоставление букв и кода}
\end{figure}

Эти коды обладают беспрефиксностью: ни один код ни одной буквы не может быть префиксом другого кода. Таким образом можно кодировать целые строки, построив дерево и пробегаясь по ноликам и единичкам. Таким образом, сочетание алгоритмов LZW и Хаффмана даёт нам zip. 

Но алгоритм Хаффмана не всегда приводит к сжатию, но иногда и к росту. Это одна из проблем. А вторая проблема так это то, что нам надо таскать с собой словарь или дерево и без них мы ничего не сможем сделать. Здесь мы сталкиваемся с проблемой под названием сериализация. \textit{Сериализация — процесс перевода структуры данных в последовательность байтов}. Тогда давайте уберем числа из дерева и сериализуем его. У нас осталось бинарное дерево. Возьмём последовательность 0111000101. Мне потребуется 8 бит, чтобы описать в каком порядке у меня встречался мой алфавит при обходе дерева: acdb (сначала влево, потом вправо). А теперь мне нужно описать структуру дерева, чтобы я мог его восстановить. Мне нужно каким-то образом написать последовательность байтиков, которые будут описывать структуру дерева, чтобы я мог прочитать строку и построить дерево. Поэтому я буду использовать три буковки: L (влево), U (вверх), R (вправо). Давайте опишем это дерево: LURLLURUURUU. В этом дереве можно убрать часть подъёма до корня, ведь если дерево кончилось, то и строка соответсвенно тоже. Остаётся LURLLURUUR. Второе предложение: вместо L и R использовать буковку D. Теперь у нас DUDDDUDUUD. Это нам даст экономию в 2 раза. 

Теперь восстановим дерево по этой записи (см. картинку снизу). Встретилась буква D, идём влево, потому что левого сына еще нет. U --- идём наверх. Опять D --- идём в правого сына, потому что левый уже есть. Снова D --- идём влево, ведь левого сына нет. Затем D --- опять влево, ведь нет левого сына. U --- наверх. D --- правый сын, ведь левый уже есть. 2 U --- Ддва раза наверх, а в конце последняя D --- создаём правого сына, ведь левый уже есть. Это и есть особенность дерева в алгоритме Хаффмана: Каждая вершина либо лист, либо имеет два ребенка.

\newpage
\begin{figure}[h]
\centering
\includegraphics[width=0.7\linewidth]{aaapic20.png}
\caption{Построение дерева}
\end{figure}

Но это запись по-прежнему избыточна. Давайте попробуем эту логику применить к кодированию. Мы пошли вниз из корня. Потом сходили по U. Но мы знаем, что мы сходили из левого сына, и сейчас деваться некуда, поэтому мы идём вправо, потому что укаждого узла должно быть два ребёнка, значит, писать D мы не будем. Потом пишем два раза D, чтобы дойти до листа. Поднимаемся два раза вверх. Идём вверх до тех пор, пока не найдём вершину, у которого нет правого ребёнка, и идём в него. Запись закончена: DUDDUU. 

А теперь восстановим по ней! Теперь U у нас имеет следующее значение: поднимайся вверх до тех пор, пока не обнаружишь вершину, у которой нет правого ребенок, и иди вправо. Пишем D, идём влево. Далее у нас U --- поднялись наверх, обнаружили вершину, у которой нет правого ребенка, и иду вправо. Дальше я встречаю D --- левый ребенок. Потом снова D --- опять левый ребёнок. Потом идём вверх пока не встречаем вершинку, у которого нет правого ребёнка, и идём вправо. Потом снова идём наверх пока не видим еще одну вершинку, у которого нет правого ребёнка, и снова идём вправо. Дерево построено. 

Количество узлов в таком дереве: 2n - 1, где n --- размер алфавита. А сколько нам потребуется для описания дерева из 2n - 1 узла? Каждый раз когда мы пишем одну букву у нас создаётся одна вершинка. И есть у нас 7 вершин и одна уже есть (корень), то нужно ровно 6 букв. Поэтому для описания структуры дерева таким образом вам понадобиться 2n - 2 бита. Значит, нам понадобиться 8 бит, чтобы знать размер алфавита, потом $N$ $\times$ 8 для хранения конкретных байт и затем нам нужно 2n - 2 бита для описания дерева. Таким образом, по описанию мы воссоздадим дерево, а далее по нему расшифруем надпись. А также можно сэкономить биты, если не хранить размер алфавита, а просто в конце описания добавить U, что будет свидетельствовать о том, что дерево уже закончилось и правых сыновей добавлять наверх уже некуда. Примерно так архиваторы и делаются.

\subsection*{Преобразование Барроуза --- Уилера}
\addcontentsline{toc}{subsection}{Преобразование Барроуза --- Уилера}

Оно очень похоже на суффиксный массив, только вместо суффиксов там делается сортировка циклических строк, ну мы это умеем делать за Nlog$^2$N, а потом применяется какой-нибудь алгоритм сжатия по последнему столбцу. 

\subsection*{Исправляющий код Хэмминга}
\addcontentsline{toc}{subsection}{Исправляющий код Хэмминга}

При передаче информации у нас возникают ошибки, и самая важная для нас ошибка --- это ошибка замены. То есть 0 на 1 и 1 на 0. Например, чтобы стабилизировать сигнал из космоса, используют троечки, то есть 3 раза отправляют одно и то же. Если у одного возникает ошибка, то у нас остаётся еще два. Но это потребляет чересчур много лишней энергии. 

У нас есть следующая последовательность:

\[
0100010000111101
\]

Нам надо научиться в ней исправлять одну ошибку. Чтобы закодировать 17 возможных вариантов (либо ошибки нет, либо в одном из 16 бит), нам потребуется не меньше 5 бит (ибо 2$^4$=16, а 2$^5$=32). Поступим мы следующим образом: мы зафиксируем 5 бит дополнительных, 16+5=21, поэтому мы создадим 21 позицию, начиная с единицы. Еще сложность в том, что в этих контрольных битах тоже может быть ошибка, и их просто так приделать нельзя. Поэтому это нужно учесть и кодировать не 17 вариантов, а 22, но к счастью для этого хватает пяти бит. 

\begin{figure}[h]
\centering
\includegraphics[width=1\linewidth]{aaapic21.png}
\end{figure}

Мы пропустили позиции, которые равны степеням двойки. Теперь вписываем в них контрольные боты. Начнём с числа 1. Пусть сумма битов на нечётных местах будет нечётной, то есть мы идём через один элемент. Для двойки: я буду считать отрезки длиной 2 с промежутком два. И суммирую то, что подчёркнуто длинной чёрточкой длины 2. Подчёркнутые определённой длины палочкой элементы в сумме должны давать нечётное число, поэтому на вторую позицию мы пишем 2. То же самое с четверкой, восьмёркой и т.д. Если в сумме получилось нечётное число, то на позицию вписываем 0, если чётное, то 1.

Теперь сломаем 20-ый бит. Сначала я проверяю те элементы, которые подчёркнуты одной палочкой. Их сумма нечётная, значит, всё хорошо. Потом проверяем элементы, которые подчёркнуты палочкой, у которой длина 2, и идут они через два. Верно. Дальше для четырёх через четыре. Тоже сумма элементов нечётная, значит, верно. Далее для 8. Тоже всё верно. Далее для 16. Нет! Находим тот самый 20-ый бит, потому что это единственный элемент, который не попал под единицу, двойку, четверку и восьмёрку, а мы ведь их уже проверили, и они оказались правильными.

\begin{figure}[h]
\centering
\includegraphics[width=1\linewidth]{aaapic22.png}
\caption{Стало (раньше под 20 была единичка, но мы исправили ее, потому что нашли ошибку)}
\end{figure}

Здесь нужно обратить внимание на уникальный набор чёрточек под каждой позицией. Мы используем N+k битов, где N --- кол-во элементов в строке и k --- исправляющие биты, таким образом, что log(N+k+1) $\leq$ k.

\newpage
\section*{Лекция 4: P и NP}
\addcontentsline{toc}{section}{Лекция 4: P и NP}

\subsection*{Определения и примеры}
\addcontentsline{toc}{subsection}{Определения и примеры}

Есть два интересующих нас класса алгоритмов: P (полиномиальные) и NP (не совсем полиномиальные, но это не точно)\footnote{В Международной лаборатории теоретической информатики в НИУ ВШЭ прошёл семинар на тему <<P = NP с точностью до переиспользования>>. Посмотреть можно по ссылке: https://cs.hse.ru/big-data/tcs-lab/news/1064461434.html}. Мы будем говорить не просто об алгоритмах, а о так называемом "языке". \textit{Язык --- это последовательность каких-то символов, в общем случае любых, но в частности в компьютере это могут быть нолики и единички.} Вот у нас есть набор всевозможных ноликов и единичек, языком L называется некоторое подмножество, возможно, бесконечного размера для этого языка. Слова принимаются, если они входят во множество слов языка, то слова принимаются алгоритмом, распознающим этот язык. Наша задача проверить, содержится ли этот набор ноликов и единичек в нашем массиве. Ответ может быть только один среди <<да>> или <<нет>>. 

\paragraph{Язык отсортированных массивов.} У нас есть все массивы, а алгоритм принимает только те из них, которые отсортированы по возрастанию. Массивы описываются двоичной последовательностью чисел.

\paragraph{Язык связных графов.} Наш язык из множества всех существующих графов принимает только связные.

\paragraph{Язык подстрок заданной строки.} Задана строка как константная и вам нужно проверить содержиться ли слова как подстрока в заданной строке.
\\

Это были примеры. Введем некоторые обозначения:\\
T(p, x) --- это сколько времени работает программа p на входе x, причем время измерено в каких-то элементарных операциях; S(p, x) --- количество ячеек памяти, которые нам нужно, чтобы программа отработала на входе х; DTIME(f(n)) --- это такой класс языков, для которого существует программа, такая что T(p, x) = O(f(n)) $\forall$ x $\in$ L: |x| = n, то есть это такой класс языков, который работает за заданное время. 

Чтобы как-то оценивать сложность, нам нужен исполнитель алгоритмов. 

\subsection*{Машина Тьюринга}
\addcontentsline{toc}{subsection}{Машина Тьюринга}

Машина Тьюринга --- это бесконечная лента, в каждой ячейке которой может быть записано слово из ограниченного алфавита. И вот у нас есть машина, она умеет кататься влево и вправо, умеет читать и писать под собой и иметь ограниченное количество состояний.

Пусть у нас на ленте записана какая-то строка, и наша машина должна проверить, является данная строка палиндромом или нет. Давайте машина будет читать символ, запоминать его, двигаться вправо до упора, сдвигаться на один влево, читать символ. Если он совпадает, то она его удаляет, а если нет, то она сразу говорит, что это не палиндром. 

\begin{figure}[h]
\centering
\includegraphics[width=0.8\linewidth]{aaapic23.png}
\caption{Машина едет...}
\end{figure}

За какую сложность мы решили эту задачу? За квадратичную. Самое главное, что она полиномиальная. Существует теория, что любой алгоритм можно реализовать с помощью машины Тьюринга. 

\subsection*{P}
\addcontentsline{toc}{subsection}{P}

Это значит, что язык, который разрешается машиной Тьюринга или другим исполнителем, исполняется за полиномиальное время. То есть

\[
P = \bigcup DTIME(p(n))
\]

Утверждается, что язык L$_1$ \textit{сводится по Карпу} к языку L$_2$, если, во-первых, L$_1$ $\leq$ L$_2$, а во-вторых, x $\in$ L$_1$ $\Longleftrightarrow$ f(x) $\in$ L$_2$. А также объединение, пересечение и дополнение полиномиальных алгоритмов тоже полиномиально. Замыкание к линии тоже, это значит, что у нас текст L* можно разбить на слова, каждое из которых принадлежит языку L. Сначала проверяем, что каждое слово принадлежит языку L за полиномиальное время, а потом проверяем, что до него тоже можно было нарезать на слова. Делаем мы это динамикой. Ею мы для каждой позиции перебираем всевозможные слова, заканчивающиеся в этой позиции. 

Детерминированная машина Тьюринга умеет находиться в нескольких состояниях, то есть она раздваивается и растраивается и каждая из них становится отдельной машиной Тьюринга. Чтобы было понятнее, представьте себе процессор с бесконечным количеством ядер и у вас условно начинает запускаться две функции перебора, и они начинают работать параллельно. 

\subsection*{NP и $\textstyle\sum_1$}
\addcontentsline{toc}{subsection}{NP и $\textstyle\sum_1$}

Это класс языков, который распознаётся недетерминированной машиной Тьюринга. P полностью входит в NP. Что такое класс $\textstyle\sum_1$? Это полиномиальная верификация сертификата с помощью детерминированной машины Тьюринга. То есть условно мы написали алгоритм сортировки, который сортирует массив. Можем ли мы за полиномиальное время проверить, отсортирован ли массив? Да, можем. Так, у нас предъявляется сертификат (сертификат --- это то, что выдаёт; некий ответ), и если мы этот сертификат можем проверить, то мы говорим, что этот язык относится к классу $\textstyle\sum_1$.

Утверждается, что $\textstyle\sum_1$ и NP это одно и то же. Доказательство: \\
$\Rightarrow$) Почему $\textstyle\sum_1$ --- подмножество NP? У нас алгоритм принадлежит $\sum_1$, потому что мы за полиномиальное время проверять сертификат. Мы можем перебором угадать ответ с помощью детерминированной машины Тьюринга достаточно эффективно. И соответственно мы пишем программу, которая перебирает все ответы на детерминированной машины Тьюринга за полиномиальное время и потом каждый из них проверяет, уже не важно на какой машине. \\
$\Leftarrow$) У нас каждый раз когда машина раздваивается, мы можем запомнить, в какую из ветвей мы пошли. Допустим, задача SAT: мы можем запомнить путь, когда формула была выполнена (то есть пусть между ложью и истиной). Таких ветвлений полиномиальное количество. И когда мы будем проверять сертификат, мы просто пройдёмся по тому же пути и проверим, действительно ли это истина. 

Следовательно, $\textstyle\sum_1$ и NP это одно и то же. Значит, определение $\textstyle\sum_1$ также является определением NP. Как решить задачу класса NP, если от нас это требуют? NP-трудный --- если любая задача из NP сводится к этому языку. NP-полный --- если любая задача из NP сводится к этому языку, а также он лежит внутри класса NP (см. картинку).

\begin{figure}[h]
\centering
\includegraphics[width=0.5\linewidth]{aaapic24.png}
\end{figure}

Нам нужен какой-то NP-полный язык. Мы докажем, что классические задачи сводятся друг к другу. Этот язык называется BH$_{1N}$. В этот язык входят тройки такие как <m, x, 1$^t$(это кол-во единичек, равное t)>, что недетерминированная машина Тьюринга возвращает единицу на словах х за время $\leq$ t. Мы описали все недетеминированные машины Тьюринга, всевозможные вводы и задали ограничение по времени. Почему этот язык относится к классу NP? Мы запускаем машину, и она по определению относится к классу NP. То есть мы сделали язык, где есть все недетеминированные машины Тьюринга. А чтобы доказать, что наша задача относится к классу NP и поэтому мы не можем её решить, а не потому, что мы криворукие, нам нужно свести её к какой-нибудь другой задаче, которая относится к классу NP. 

Давайте мы сведем язык BH$_{1N}$ к задаче SAT --- это задача проверки булевой функции. Что нам понадобиться? У нас есть описание НМТ. В каком-то из состояний она либо не приходит к успеху, либо приходит. Сведение BH$_{1N}$ к задаче SAT заключается в том, чтобы каждое это раздваивание описать в виде переменной в логической функции. А формулой логической проверяется, привёл ли этот выбор раздваивание к успеху или нет. Наши операции $\lor$, $\land$ и $\lnot$ --- это базис всех логических операций в компьютере и булевая функция как раз таки их описывает. 

CNF --- конъюктивно-нормальная форма. Как доказать, что задача CNFSAT NP-полная? Нужно SAT свести к CNFSAT. 3-SAT --- это конъюкция из трёх переменных. Нужно CNFSAT свести к 3-SAT. 

\paragraph{Задачи о поиске клики и независимого подмножества.} Они друг к другу сводятся: если мы хотим решить задачу о клике, а умеем решать задачу о независимых множествах, то нужно мы все рёбра инвертируем. А как нам доказать, что задача о поиске НЗМ также является NP-полной? Нам надо какую-то из наших задач, про которые мы уже знаем, что она NP-полная, свести к задаче о поиске независимого множества. И сводить мы будем 3-SAT. Класс NP совпадает с классом $\sum_1$, поэтому мы можем предъявить на детерминированной машине Тьюринга алгоритм проверки сертификата, сказать, что задача о поиске НЗМ сводится к $\sum_1$, а это и есть NP. А любую задачу из класса NP можно свести к задаче о поиске НЗМ заданного размера. 

Пусть у нас есть данный 3-SAT:

\[
(x_1|x_2|!x_4) \& (x_1|!x_2|x_4)
\]

$\Rightarrow$) Мы сведём по три вершины на каждую конъюкцию. У нас должно быть ребро внутри каждого такого дизъюнкта, то есть треугольничком, а также каждое ребро связываем с его отрицанием. И теперь на получившемся графе мы хотим найти НЗМ размера k, где k --- количество дизъюнктов. Как оно найдётся? Очевидно, что в каждом треугольнике будет выбрана ровно одна вершина. Выберем оба x$_1$. Ответом будет <<да>> или <<нет>>. Почему это верно? В каждом треугольнике мы выбираем по одной вершинке и при этом гарантируется, что не будет выбрана переменная и её отрицание одновременно, потому между всеми такими парами будет проведено ребро. А значит, когда мы предъявим сертификат, какие вершины выбраны, мы подставим эти значения: те, что выбраны, истинны, остальные ложны. Получаем, что гарантированно получим истину. Мы доказали, что это задача класса NP. И если мы свели NP-полную задачу к частному случаю задачи, относящейся к классу NP, то эта задача NP-полной. 

$\Leftarrow$)Допустим, у нас есть какое-то решение, но не существует НЗМ размера k. Давайте все переменные, помеченные истиной возьмём в независимое множество. Предположим, что у нас не выбрано ни одной вершины. Значит, какая-то из дизъюнкт приняла значение ЛОЖЬ, а такого быть не может. Следовательно, хоть какая-то вершина да выбрана. Если одна в треугольнике, то мы радуемся, а если нет? Если больше одной, то выбрасываем все, кроме одной. Вуа-ля, независимое множество размера k. 

\begin{figure}[h]
\centering
\includegraphics[width=0.5\linewidth]{aaapic25.png}
\end{figure}

\newpage
\section*{Лекция 5: Оптимизация перебора}
\addcontentsline{toc}{section}{Лекция 5: Оптимизация перебора}

\subsection*{Градиентный спуск}
\addcontentsline{toc}{subsection}{Градиентный спуск}
\paragraph{Задача о холме.} Допустим, у вас есть холм вы взбираетесь на него. Ничего не видно, но у вас есть палочка. Как взобраться на вершину? Слева и справа от себя вы можете понять, где выше, и идти туда. Когда вы дойдёте до вершины, у вас возникнут вопросы: во-первых, с какой точностью вы достигли её, а во-вторых, идти вы можете очень долго. Предположим, что палочка у нас переменной длины. Сначала мы будем делать большие шаги, а потом постепенно их уменьшать, причём окажется, что нам надо уменьшать их вдвое.

Мы ищем вершину холма. Где мы должны в начале встать? Мы должны расположиться посередине. И если мы идём оттуда, то шаги должны быть в правую и левую четверти. И так продолжаем потихоньку двигаться в нужном направлении, уменьшая шаг на два. Это называется \textit{градиентный спуск}. Спуск, потому что шаги уменьшаются. Количество измерений равно количеству измерений в бинпоиске. 

\begin{figure}[h]
\centering
\includegraphics[width=1\linewidth]{aaapic26.png}
\end{figure}

Теперь в двумерном пространстве. Почему бы нам не потыкать точки и по х, и по у и пойти в ту из них, которая выше всего. Выбрали максимум из четырёх направлений и пошли туда. Теперь давайте посмотрим, какие шаги нам делать, чтобы достигнуть абсолютно любой точки нашего холма. И так, мы стоим посередине, и первый шаг должен быть равен половине стороны квадрата, чтобы мы могли попасть на сторону квадрата. Таким образом, уменьшая шаг вдвое, сумма шагов будет равна стороне квадрата.

\newpage
\begin{figure}[h]
\centering
\includegraphics[width=0.7\linewidth]{aaapic27.png}
\end{figure}

\subsection*{Метод Монте-Карло и квадратики}
\addcontentsline{toc}{subsection}{Метод Монте-Карло}

\paragraph{Задача об участке земли.} Эта задача не имеет никакого математического решения. Вы имеете участок, и вам нужно выбрать самое высокое место, чтобы построить ветряк. Вы купили квадракоптер, чтобы измерять высоту. Вы можете произвести N измерений. Один из способов --- это померить в случайных точках и определить самую высокую как вершину вашего участка. По-честному делаем N измерений. Также можно нарезать весь участок на квадратики и в каждом квадратике произвести измерения. Но метод Монте-Карло лучше, потому что если нарезать на квадраты, то мы воспользуемся не всеми измерениями, а меньше. Допустим, мы нарезали квадратики достаточно мелко, чтобы там внутри функцию можно было принять за выпуклую. То есть часть измерений мы потратили, чтобы выявить общую картину, сами квадратики, а оставшимися измерениями мы возьмём перспективные квадратики (те, где измерения достаточно большое) и улучшим в них, сделав пару шагов нашего градиентного спуска. 

\paragraph{Задача о распределении студентов.} У нас есть 200 студентов. Для каждого есть результаты тестирования по математике, программированию и пол по паспорту. Требуется их распределить. Нужно распределить их по среднему значению M$_i$, P$_i$ и G$_i$. Употреблять термин <<дисперсия>> мы пока не будем. У нас есть прекрасная функция shuffle(), которая перемешивает последовательности. Сначала написать 25 единичек, потом 25 двоек и 25 восьмёрок. Далее случайным образом перемешать. И это нам гарантирует, что в каждой группе ровно 25 человек. Как нам ещё улучшить распределение? Нужно написать функцию оценки. Допустим, она уже есть, то есть у нас написано сколько стоит отклонение от среднего по математике, программированию и т.д. А дальше оно либо квадратичное, либо линейное, то есть модуль разности суммируется по всем группам (или квадрат модуля разности). В общем, функцию за нас уже написали и мы можем её реализовать. И теперь мы пытаемся насчитать ту функцию штрафа за отклонение и, допустим, мы успели ее написать. Применяем, получаем какое-то отклонение. 

Но, может быть, поскольку наш метод Монде-Карло никогда не попадёт во что-то хорошее, ведь область определения довольно большая. И тут мы должны сделать аналогию с нашим градиентным спуском. Мы должны для нашего решения посмотреть соседей и среди всех соседей выбрать лучшего. Как мы это сделаем? У нас какая-то расстановка, распределение. Допустим, по 2 человека на 3 группы. Мы пробуем переставить двух людей и смотрим, а не стало ли лучше. И тут мы действуем двумя способами: либо перебрать всевозможные пары людей и выбрать наибольшее улучшение, тогда это будет градиент, либо первая попавшаяся пара, которая даёт улучшение, сразу применяется и уже мы делаем следующее. Какой из методов лучше зависит от времени. Менять местами людей гораздо быстрее, чем писать заново.

В какой-то момент обмен любой пары людей перестанет приносить улучшения. И вот мы нашли локальный максимум, а дальше снова генерируем перестановки и ищем его уже там. И так делаем пока не надоест.

Существует несколько алгоритмов основанных на этой идее. Проблема заключается в том, что мы упираемся в локальный максимум. Например, алгоритм имитации отжига. 

\paragraph{Задача о нефтепроводе.} У нас есть трубы и насосные станции в виде графа. Мы можем у насосной станции поставить датчик, который измеряет, сколько у нас протекает через все трубы. И нам надо проконтролировать абсолютно все рёбра графа. Это задача о минимально контролирующем множестве, то есть нам нужно выделить не более k вершин таким образом, чтобы у каждого ребра хотя бы один его конец был помечен. Введём ограничения: N (рёбер) $\leq$ 2000, k $\leq$ 40. 

Жадный алгоритм здесь не работает. Идея решения: давайте рандомно пометим k точек и посмотрим, что получилось. Нам нужно придумать функцию оценки качества неправильного решения. Если мы нашли хоть какое-то покрытие, то мы уже молодцы, дальше можно не думать. Теперь нам нужно определиться, какое из неправильных решений более правильное: чем больше труб покрыто, тем лучше. Если вы в терминах штрафа хотите мыслить, то количество непокрытых. Если у вас есть время, то генерируем функцию оценки, пока не выйдет TL. А как собственно функцию оценки сделать? Если у нас 2000 вершин, то у нас может быть порядка 2 000 000 рёбер. И если вы каждое ребро будете отдельно проверять, то у вас будет не так много решений. Давайте быстро научимся считать качество нашего решения. То есть мы выбрали k вершин и нам нужно понять, сколько рёбер покрыто нашими k вершинами. Поэтому мы воспользуемся суммой степеней помеченных вершин. У нас их мало: всего k. И если мы их просуммируем, то это будет работать сильно быстрее, за O(k). Но теперь у нас еще одна проблема: мы одно и то же ребро можем посчитать дважды. Нам нужно только перебрать все пары помеченных вершин и проверить, есть ли между ними ребро. Если есть, то из счётчика вычесть единичку. Делается это всё за O(k$^2$). 

Теперь задаём локальный оптимум. Мы будем переносить одну, две или три пометочки и смотреть, стало ли лучше. Нам нужно будет вернуть ребро, потому что оно было вычтено как помеченное с двух сторон, а теперь оно помечено с одной стороны и возможно какие-то рёбра возможно выкинуть, если они стали помеченные с двух сторон. Но это можно сделать за O(2k). То есть все эти махинации мы делаем только для свежих вершин. Это самое правильное решение с всероса с одним небольшим дополнением. Если у вас есть одна длинная цепочка, то это дело можно быстро ускорить. Каким образом она должна выглядеть? Нам нужно через одну выбирать вершины. И такая же идея, если у нас есть длинная цепочка, которая ведёт не в лист, а просто длинная цепочка вершин. Это здорово ускоряет алгоритм. 

\subsection*{Итоги лекции}
\addcontentsline{toc}{subsection}{Итоги лекции}

Я предлагаю вам схему, которая будет довольно универсально работать. Сначала вы пишите генерацию рандома, убеждаетесь, что она работает. Затем вы добавляете функцию оценки и можете применить метод Монте-Карло. То есть генерируете рандомно и выбираете лучший. Также можете написать еще более быструю функцию оценки, если думаете, что у вас очень медленно считается. И после этого вы запускаете это на всё время, что у вас есть, и надеетесь, что всё будет хорошо. Какие есть подводные камни в этом всём? В части домашек (например, задача про поиск Гамильтонова пути) вы рандом никогда не накидаете даже близко к жаднику. В таком случае нужно запустить жадник и оптимизировать его. И в реальности это даст результат лучше, чем рандом. 

\newpage
\section*{Лекция 6: Персистентные структуры данных}
\addcontentsline{toc}{section}{Лекция 6: Персистентные структуры данных} 

\textit{Персистентные структуры данных} --- это структура данных, которая позволяет как минимум получать свои предыдущие версии. Мы хотим добиться того, чтобы прочитать предыдущие версии, а может быть, даже унаследоваться от них. В реальности это широко распространено. Какие бывают персистентности?

\subsection*{Виды и примеры}
\addcontentsline{toc}{subsection}{Виды и примеры}

\begin{enumerate}
\item Частичная --- читать вы можете любую версию, а писать только последнюю.
\item Полная --- вы можете любую версию писать.
\item Конфлюэнтная --- позволяет сливать между собой структуры данных.
\item Функциональная --- ваши переменные присваиваются только при их создании.
\end{enumerate}

Нам интересны пока только первые две. Возьмём обычный массив с буковками. Самый простой способ организовать персистентность --- это при каждом изменении делать полную копию того, что было раньше. Это очень накладно, но тем не менее такой способ существует. Существуют и более продвинутые способы. Например, копирование пути: мы делаем изменение только в том пути, в котором это изменение произошло. 

Пусть у нас есть бинарное дерево поиска. И мы хотим добавить сюда один элемент, сохранив персистентность. (см. картинку снизу). Добавлять его мы будем точно так же, как и раньше, но, поскольку нам понадобится и старая, и новая версия, то всё, что было на пути, мы будем копировать. Левую ветку из старого дерева мы полностью копируем, потому что в нём никаких изменений не произошло, а в правом дереве мы обязательно создаём новый элемент, в нашем случае это с`. Чтобы дойти до h, нам нужно свернуть влево, а значит, правое поддерево мы полностью копируем. А дальше приписываем h. Всё! Теперь мы можем осуществлять поиск и в первом, и во втором дереве. Сложность происходящего в худшем случае --- логарифм узлов, если дерево было сбалансированным. Если нет, то O(N). Если дерево было сбалансированным, то на хранение новой версии дерева уходит логарифм памяти. Да и времени тоже, но константа хуже, чем обычно. Персистентность это полная. 

\begin{figure}[h]
\centering
\includegraphics[width=0.7\linewidth]{aaapic28.png}
\end{figure}

Это всё очень круто, до тех пор пока у нас не появляется ссылка на предка. В момент, когда у нас появляется ссылка на предка, у каждого элемента у нас происходит проблема. В одной версии у нас может быть один предок, в другой версии --- другой. Как это хранить? Непонятно. Вернёмся к этому чуть-чуть попозже.

Теперь реализуем персистентный стек. Давайте реализуем наш стек в виде односвязного списка. Предположим, наш стек состоял из элементов 1, 3 и 5. Как у нас устроен неперсистентный стек? Мы добавляем, условно, цифру 7, от него протягиваем ссылку на предыдущий элемент и устанавливаем, что стек у нас заканчивается на 7. Чтобы сделать персистентный стек, мы заведём массив версий. Допустим, в первой версии у нас 1. Мы добавляем к единичке троечку. Мы создаём новый элемент, поскольку это push, делаем на него ссылку на предыдущий элемент и кидаем её в новую версию (см. картинку сверху). Теперь добавим пятёрку: она попадает в ячейку с новой версией, а её ссылка начинает показываться на предыдущий элемент списка.

А если я хочу четвёрку добавить ко второй версии? Всё то же самое, но только теперь четвёрка у нас указывает на тройку. 

\begin{figure}[h]
\centering
\includegraphics[width=0.5\linewidth]{aaapic29.png}
\end{figure}

Давайте сделаем pop из четвёртой версии. Пятая версия будет указывать на тройку, потому что перед ней стоит четвёрка, которую мы удалили. 

\begin{figure}[h]
\centering
\includegraphics[width=0.5\linewidth]{aaapic39.png}
\end{figure}

Получается такое корневое дерево, которое реализовывает полную персистентность. Сложность поддержки и по времени, и по памяти --- O(1). 

Вернёмся к дереву поиска. В прошлом семестре наше неявное декартово дерево имело merge и split. Можно ли декартово дерево сделать персистентным? Конечно, по такой же схеме. Ссылки на предков будут всё портить, и с этим ничего не поделаешь. Разберём новый пример. 

Дерево отрезков. Мы хранили его как массив элементов. Мы пока не умеем делать персистентный массив. Я предлагаю такое решение: вместо палочек в дереве отрезков нарисовать стрелочки. Что означают стрелочки? Они значат, что теперь мы будем хранить дерево отрезков не на массиве, а в виде обычных нод как в бинарном дереве: есть левая нода, есть правая нода. Всё то же самое, что и в бинарном дереве. То есть мы будем реализовывать бинарное дерево не снизу, а сверху. Это хуже по памяти, но позволит нам добавить персистентности. А можно ли совместить персистентность с проталкиванием обещания? В принципе, да. Смотрите, когда мы обещание проталкиваем, мы по-прежнему совершаем логарифм действий. Мы идём по какому-то пути, и если наш узел что-то обещал, то мы проталкиваем его непосредственных детей. Как это реализовать в случае, если у нас есть персистентность? В непоследней версии мы ничего менять не будем, а в изменённой версии, то есть если мы протолкнули какое-то обещание, мы его проталкиваем только в новой версии этого дерева. Но в старой версии вы ни в коем случае обещания не проталкиваете. 

В любой деревовидной структуре понятно: либо в явном виде, либо в односвязных списках, которые будут ветвиться. А также важным критерием является отсутствие ссылки на предка. 

Для начала реализуем персистентный двусвязный список. Нам он дан и наша задача реализовать частичную персистентность. Что такое изменение в двусвязных списках? Добавить ноду, удалить и изменить ключ. Допустим, мы хотим изменить ключ в ноде. Что нам потребуется хранить? Сначала разберём реализацию массива.

Вы хотите реализовать частичную персистентность обычного массива. Нам потребуется хранить в каждом элементе все изменения, которые в массиве производились, то есть все предыдущие значения в этом элементе. На самом деле каждый элемент в массиве представляет собой массив пар, состоящий из версии и её значения. Он естественным образом получится упорядоченным по версиям. Предположим, мы хотим найти вторую версию. Мы бинпоиском ищем в массиве пар максимальное значение версии, не превосходящее то, которое мы спрашиваем. Но это не совсем универсальный способ. Поэтому мы осуществляем полное копирование массива, когда наполнилось до какого-то определённого предела. Допустим, берём по 1000, а нам нужна 10500-ая версия, получается, это 10-ая копия. А далее делаем вышеописанные действия. Наша главная цель --- это сохранить логарифм в сложности.

Эту же идею мы сейчас применим к персистентному двусвязному списку. Кроме ключей, которые мы будем хранить точно так же, будем также держать у себя список prev и next. И когда у нас добавляется новый элемент, мы говорим, что в списке next'ов новая версия и показываем на неё. То же самое с prev.

\begin{figure}[h]
\centering
\includegraphics[width=0.5\linewidth]{aaapic40.png}
\end{figure}

К сожалению, записи этой лекции прошлых годов не сохранились, поэтому данный конспект будет дополнен сразу после проведения занятия для нашего года.


\newpage
\section*{Лекция 7: Паросочетания}
\addcontentsline{toc}{section}{Лекция 7: Паросочетания} 

\subsection*{Определения}
\addcontentsline{toc}{subsection}{Определения}

\textit{Двудольный граф --- это граф, вершины которого можно разделить на две непересекающиеся группы (доли) так, что каждое ребро соединяет вершину из одной доли с вершиной из другой доли. Другими словами, внутри каждой из долей не должно быть ребер. Паросочетание --- это набор попарно несмежных рёбер.} Нам нужно построить максимальное паросочетание. Нам нужно ввести некоторые определения. \textit{Чередующаяся цепь --- это путь, в котором ребра попеременно принадлежат и не принадлежат паросочетанию. Дополняющая цепь --- это чередующаяся цепь, оба конца которой свободны, то есть не покрытые рёбрами из паросочетания.} С помощью дополняющей цепи можно увеличить мощность паросочетания на 1. Это позволит нам решить задачу о поиске максимального паросочетания. 

\subsection*{Задача о поиске максимального паросочетания}
\addcontentsline{toc}{subsection}{Задача о поиске максимального паросочетания}

Воспользуемся жадным алгоритмом: мы ищем дополняющие цепи. Мы берём какое-то ребро (см. картинку)

\begin{figure}[h]
\centering
\includegraphics[width=1\linewidth]{aaapic41.png}
\end{figure}

Можем ли мы увеличить паросочетание до дополняющей цепи? Ну, нет, не можем. Что мы делаем? Берём это ребро и принимаем в паросочетание.

\newpage
\begin{figure}[h]
\centering
\includegraphics[width=1\linewidth]{aaapic42.png}
\end{figure}

Берём следующую дополняющую цепь. 

\begin{figure}[h]
\centering
\includegraphics[width=1\linewidth]{aaapic43.png}
\end{figure}

Следующим шагом мы смотрим, можно ли еще увеличить цепь с двух сторон.

\begin{figure}[h]
\centering
\includegraphics[width=1\linewidth]{aaapic44.png}
\end{figure}

Поэтому мы красим ранее выбранное ребро в красный, поскольку мы добавили рёбер,

\newpage
\begin{figure}[h]
\centering
\includegraphics[width=1\linewidth]{aaapic45.png}
\end{figure}

а затем сразу же перекрасить, так как мы с вами обсуждали: чтобы увеличить размер паросочетания. 

\begin{figure}[h]
\centering
\includegraphics[width=1\linewidth]{aaapic46.png}
\end{figure}

То есть мы ищем какую-то цепь максимальную, и если вдруг оказалось, что она нечётной длины, то крайние рёбра принимаем в паросочетание. 

Чтобы доказать, что это работает, мы должны доказать следующее утверждение. 

\subsection*{Лемма Бержа}
\addcontentsline{toc}{subsection}{Лемма Бержа}

Паросочетание максимально тогда и только тогда, когда в графе не существует дополняющей цепи.

$\Rightarrow$) Если есть дополняющая цепь, то мы увеличили паросочетание, а значит, до этого оно не было максимальным.

$\Leftarrow$) Рассмотрим M $\oplus$ M'

\newpage
\begin{figure}[h]
\centering
\includegraphics[width=1\linewidth]{aaapic47.png}
\end{figure}

Что мы можем про него сказать? 

\begin{figure}[h]
\centering
\includegraphics[width=0.6\linewidth]{aaapic48.png}
\end{figure}

Во-первых, абсолютно у любой вершины её степень будет не больше чем 2. Это происходит потому, что у нас все рёбра, и красные, и синие, взяты из какого-либо паросочетания. А значит, если мы этот граф разобьём на компоненты, то это будет либо путь, либо цикл. Понятно, что в таком графе остаются только цепи и циклы чётной длины. В цикле чётной длины цвета рёбер будут чередоваться, ведь в паросочетании запрещено проводить два ребра из одной вершины. И кроме циклов обязательно найдётся хотя бы один путь, состоящий из нечётного количества рёбер, при чём красных рёбер у него больше быть не может. У нас есть какой-то путь, причём синих там больше, чем красных. Это значит, что в пути мы можем синие рёбра перекрасить в красный и получить дополняющую цепь, которой у нас по определению нет. Пришли к противоречию. 

\subsection*{Алгоритм Куна}
\addcontentsline{toc}{subsection}{Алгоритм Куна}

Используя DFS, мы будем искать самые длинные пути, которые нельзя увеличить. Как это писать? Мы создадим вершину S и соединим её со всеми вершинами левой доли. А также вершину T, которую соединим с вершинами из правой доли. Идея: рёбра, которые у нас не входят в паросочетание, мы будем ориентировать в сторону правой доли. А далее запускаем поиск от S до T по всем этим рёбрам. Когда у нас какой-нибудь путь нашёлся, берём в паросочетание, красим в красный цвет и разворачиваем.

\begin{figure}[h]
\centering
\includegraphics[width=0.3\linewidth]{aaapic49.png}
\end{figure}

По этому ребру мы пройти больше не можем. А на следующем шаге мы ищем уже другой путь. Он нам подходит, поэтому мы все рёбра переориентируем в обратную сторону.

\begin{figure}[h]
\centering
\includegraphics[width=0.3\linewidth]{aaapic50.png}
\end{figure}

У нас получилась та самая увеличивающаяся цепочка. И так далее, пока мы не переориентируем все вершины, лежащие между рёбрами. Каждый раз мы DFS запускаем заново для ещё непереориентированных рёбер. Сложность такого алгоритма --- О(NM), где N --- кол-во вершин и M --- кол-во рёбер. 

Но можно решить эту задачу более житейским способом и даже быстрее. Допустим, у нас есть работники и нам нужно сопоставить им работу. 

\begin{figure}[h]
\centering
\includegraphics[width=0.5\linewidth]{aaapic51.png}
\end{figure}

Начнём с первого. Он просматривает список всех работ, ту, которая ещё никем не занята, и занимает её. Второй работник делает то же самое. 

\begin{figure}[h]
\centering
\includegraphics[width=0.5\linewidth]{aaapic52.png}
\end{figure}

Дошли до третьего работника: ему придётся выгнать кого-нибудь с его работы, чтобы получить себе место. Что он делает? Он накрывается домиком и говорит: <<Товарищ 1, ищите-ка вы себе новую работу.>> Он отвечает <<Ладно>> и тоже накрывается домиком и забирает работу у второго товарища. Говорит ему, чтобы он заново искал себе работу. После этого второй ищет себе новую работу, и все радуются и убирают шапочки. 

\begin{figure}[h]
\centering
\includegraphics[width=0.5\linewidth]{aaapic53.png}
\end{figure}

Идём к четвёртому. У него нет работы, поэтому он надевает шапочку и пытается отобрать работу у второго человека. Второй снова ищет работу и уже идёт на неё.

\begin{figure}[h]
\centering
\includegraphics[width=0.5\linewidth]{aaapic54.png}
\end{figure}

Давайте добавим пятого работника. Он забирает работу у второго и надевает шапочку. Поотбираем работу у соседей и придём к выводу, что у пятого-то её нет. 
\newpage
\begin{figure}[h]
\centering
\includegraphics[width=0.5\linewidth]{aaapic55.png}
\end{figure}

Идея: в первую очередь вы должны пройтись по всем соседям. Если есть какой-то сосед, который еще пока что не занят, то мы ему сразу сопоставляем. Это довольно важная оптимизация. Если вдруг таких соседей нет, то пытаемся у кого-нибудь отобрать это место.

Теперь схематичный код:

\begin{verbatim}
kuhn(now):
    for neig in g[now]:
        if w[neig] == -1:
            w[neig] = now
            return True
    block[now] = True
    for neig in g[now]:
        if !block[w[neig]]:
            if kuhn(w[neig]):
                w[neig] = now
                return True
    return False
\end{verbatim}

Ключевые моменты: если человек получил работу, он её уже никогда не потеряет. Да, он может её сменить. Это первое, а второе --- это то, что вы весь массив блокировок сделаете фолсами перед тем, как запускать алгоритм Куна. 

\begin{verbatim}
for i = 1, n:
    block = [False]
    kuhn(i)
\end{verbatim}

Какую еще можно оптимизировать? Вы можете сделать блокировку, которая не снимается. Раньше мы каждый раз её снимали. Если он вернул false, то true он больше никогда не вернёт. А также надо всегда выбирать меньшую долю для запуска. 

\subsection*{Задача о минимальном вершинном покрытии}
\addcontentsline{toc}{subsection}{Задача о минимальном вершинном покрытии}

Это задача относится к классу NP. Как его эффективно искать на двудольном графе? Здесь следует вспомнить теорему Кёнига: максимальное паросочетание равняется минимальному вершинному покрытию. 

По аналогии с нашими рассуждениями о чередующихся цепочках. Мы наши рёбра, входящие в паросочетание, ориентируем справа налево, а рёбра, не входящие в паросочетание, --- слева направо. 

\begin{figure}[h]
\centering
\includegraphics[width=0.3\linewidth]{aaapic56.png}
\end{figure}

А далее мы для всех свободных вершин (свободные вершины --- это те, в которые не ведёт красное ребро) запустим обход. Пометим свободную и отметим все достижимые из неё. Помеченные вершины назовём L$^{+}$ и R$^{+}$, а остальные L$^{-}$ и R$^{-}$. 

\begin{figure}[h]
\centering
\includegraphics[width=0.25\linewidth]{aaapic57.png}
\end{figure}

Давайте обобщим. Какие рёбра в графе могут существовать, если нарисовать множество вершин в графе? Если мы целиком возьмём множества L$^{-}$ и R$^{+}$, то все рёбра хотя бы на одном своём конце будут иметь помеченную вершину. Почему?

\begin{figure}[h]
\centering
\includegraphics[width=0.36\linewidth]{aaapic58.png}
\caption{Пунктирные рёбра не могут существовать, а обычные --- да}
\end{figure}

Потому что рёбра из L$^{-}$ в R$^{-}$ будут помеченны в L$^{-}$, а рёбра из L$^{+}$ в R$^{+}$ будут иметь пометку в R$^{+}$, а рёбра между L$^{-}$ и R$^{+}$ будут с двух сторон помеченны. А рёбер между L$^{+}$ и R$^{-}$ вообще не существует, как мы видим из рисунка. 

Почему |L$^{-}$| + |R$^{+}$| = размеру максимального паросочетания? Что мы можем сказать об их размерах? В L$^{-}$ все вершины покрыты, там нет свободных вершин. Это значит, что все красные рёбра, которые идут из R$^{-}$ в L$^{-}$ соответствуют только одной вершинке: больше одного быть не может, это ведь паросочетание, а быть равным нулю тоже нет, потому что иначе они были бы свободные. Сколько рёбер из R$^{-}$ в L$^{-}$, столько же вершин в L$^{-}$.

Теперь предположим, что мы пришли в R$^{+}$ из L$^{+}$. Пусть этой вершине из R$^{+}$ не сопоставлено ни одного красного ребра. В чём проблема? Мы получили какую-то чередующуюся цепочку, которую можно увеличить, чего не может быть, ведь мы доказали теорему в начале лекции, что мы находим паросочетание в тот момент, когда у нас не остаётся дополняющих цепей. 

Выходит, что каждое ребро сопоставлено R$^{+}$ и L$^{+}$. Вот и всё доказательство. 

\subsection*{Задача о максимальном независимом множестве в двудольном графе}
\addcontentsline{toc}{subsection}{Задача о максимальном НЗМ в двудольном графе}

Утверждается, что дополнение вершинного покрытия --- это независимое множество. Поэтому мы выбираем R$^{-}$ и L$^{+}$.

\newpage
\section*{Лекция 8: Потоки}
\addcontentsline{toc}{section}{Лекция 8: Потоки} 

Разберем пример.

Возьмём граф, в котором есть исток S и сток T. При этом возле рёбер приписаны числа, которые называются пропускной способностью этого ребра. 

\subsection*{Метод Фолда --- Фалкерсона}
\addcontentsline{toc}{subsection}{Метод Фолда --- Фалкерсона}

Нам необходимо понять, какое максимальное количество ресурсов можно доставить от истока к стоку. Для каждой вершины необходимо фиксировать, сколько в неё втекает и сколько вытекает. Идея очень простая: ищем какой-нибудь путь и по нему пускаем наш поток.

\begin{figure}[h]
\centering
\includegraphics[width=0.42\linewidth]{aaapic59.png}
\end{figure}

Предположим, что наш путь --- это 7-2-4-12, и вдоль этого пути я пускаю поток. Поток у нас будет весом 2, потому что это самая маленькая из всех пропускных способностей, и запоминаю это где-то. А далее я уменьшаю веса рёбер на этом пути на 2.

А теперь наш хитрый ход! Нам нужно создать обратные рёбра с пропускной способностью 2.

\begin{figure}[h]
\centering
\includegraphics[width=0.42\linewidth]{aaapic60.png}
\end{figure}

Какой физический смысл имеют обратные рёбра с весом 2? Если у вас из вершины А в вершину Б ежесекундно проходит 2 кубометра воды, то вы можете организовать доставку из Б в А двух кубометров воды, остановив подачу воды из А в Б. 

Допустим, мы выбираем другой путь, и там тоже минимальное ребро имеет вес 2. Делаем всё то же самое и уже приплюсовываем эту двойку к той двойке, которая была в первом пути, и так далее.

Этот процесс в большинстве случаев когда-нибудь закончится. И, соответственно, мы накопим некоторую величину потоков и узнаем ответ. А за сколько закончится этот процесс? Допустим, мы знаем ответ --- |F| (величина максимального потока). Каждый раз, если у нас веса рёбер целочисленны, то мы каждый раз в поток будем пускать единичку, ведь единица --- это самый маленький вес. Один путь мы ищем за O(Е) в худшем случае. Но можно доказать, что этот алгоритм не будет работать, если в нём иррациональные веса рёбер. Попозже мы этот алгоритм оптимизируем.

Как это реализовать? Запускаем DFS, а как уменьшить вес рёбер? Для каждой вершинки мы храним список: куда она ведёт и какая у неё пропускная способность. Как у нас в алгоритмах DFS и BFS восстанавливается путь? Это просто последовательность номеров вершин. А теперь вспомним, что нам нужно рёбра уменьшить. Если мы будем пользоваться только информацией восстановленной последовательности, то как это будет происходить: из начала нужно посмотреть на всех её соседей, найти вершину, которая идёт далее по пути, и для неё уменьшить. Но это не испортит асимптотику, но портит константу и всё равно это неприятно писать. Это можно улучшить, если хранить соседей в unordered\_set. И тогда мы сможем за О(1) находить ребро. Константа всё равно плохая. Попробуем обойтись обычным вектором. Поэтому теперь мы будем запоминать не только вершину, по которой пришли, но и ребро. Значит, теперь мы храним номер вершины, откуда вы попали, и порядковый номер v$_2$ в списке соседей v$_1$. И теперь, когда вы будете восстанавливать ответ, вы знаете не только номера вершин, но и какой номер у них был в списке соседей у предыдущей вершины, и за О(1) это исправляете. 

Другая проблема --- это обратные рёбра. Заранее создадим обратные рёбра с пропускной способностью 0 прямо на этапе построения графа. Кроме того, чтобы уменьшить рёбра в пути, нам нужно ещё увеличить пропускную способность обратных рёбер. Как только мы считали из входных данных, что у нас есть какое-то ребро с пропускной способностью, мы его тут же добавляем в список соседей. И прямо сейчас мы создаём обратное ребро, а дальше храним индекс элемента в списке соседей (мы ведь добавили ребро и его обратное ребро одновременно, поэтому индекс обратного ребра --- это size-1). 

\begin{figure}[h]
\centering
\includegraphics[width=0.42\linewidth]{aaapic61.png}
\end{figure}

То есть мы сначала уменьшаем пропускную способность ребра с помощью хранения в обходе в ширину номера в списке соседей. Далее мы знаем индекс обратного ребра в списке соседей той же вершины и увеличиваем его. (см. картинку: мы сначала прыгаем в v$_2$, потом оттуда идём в список соседей v$_2$ и идём по индексу к v$_1$, а далее увеличиваем его пропускную способность). 

\subsection*{Минимальный разрез}
\addcontentsline{toc}{subsection}{Минимальный разрез}

Минимальный разрез в графе --- это минимальная суммарная чиселка, которая написана на рёбрах, такая что при удалении этих рёбер невозможно попасть от истока к стоку. Удивительным образом оказывается, что вес минимального разреза равен весу максимального потока. Как найти сами рёбра, входящие в минимальный разрез? На самом деле довольно просто ищется (см. картинку). Мы сделали на графе максимальный поток, запускаем обход в ширине по ненулевым рёбрам. Понятно, что до T мы не дойдём, потому что у нас уже нет ни одного пути. То есть есть какие-то достижимые вершины из S и из них торчат рёбра нулевого веса в недостижимые вершины. И именно их мы и будем резать, которые из достижимых торчат в нулевые (см. картинку).

\begin{figure}[h]
\centering
\includegraphics[width=0.25\linewidth]{aaapic62.png}
\end{figure}

Мы не ходим по нулевым рёбрам.

А как искать максимальный поток в графе, где рёбра неориентированы? Мы создаём прямое и обратное ребро весом неориентированного ребра. То есть если в ориентированном графе обратное ребро по умолчанию имело вес 0, то в неориентированном графе обратное ребро имеет такой же вес, что и прямое. Разрез ищете тот же самый.

\subsection*{Алгоритм Эдмондса --- Карпа}
\addcontentsline{toc}{subsection}{Алгоритм Эдмондса --- Карпа}

Если мы ищем не какой-то путь, а кратчайший по количеству рёбер путь, то это становится алгоритмом Эдмондса-Карпа. Пользуемся мы обходом в ширину, ищем поток в графе. Его сложность становится O(VE$^2$).

Чтобы доказать, что у нас всё хорошо работает, нам понадобится некоторая лемма: если у нас существовал какой-то кратчайший путь, то после всех наших действий в алгоритме он не может уменьшиться. То есть он либо останется прежним, либо увеличится.
А также лемма о критическом ребре: пусть P - кратчайший путь из S в T, назовём ребро (v, u) критическим, если его вес c(v, u) минимален среди всех весов рёбер на пути P. Тогда ребро может становиться критическим не более O(V) раз.

Каждое ребро может становиться критическим O(V) раз. Суммарно все ребра становятся критическими O(VE) раз. После пускания потока вдоль кратчайшего пути, хотя бы одно ребро становится критическим, значит всего может быть не более O(VE) увеличений потока. Поиск кратчайшего пути мы делаем с помощью обхода в ширину, следовательно, время работы алгоритма O(VE$^2$).

\newpage
\section*{Лекция 9: Базовая геометрия}
\addcontentsline{toc}{section}{Лекция 9: Базовая геометрия}  

Какие формулы нам понадобятся?
\begin{enumerate}
\item $\sqrt{x^2+y^2}$
\item $(a, b) = x_1x_2 + y_1y_2 = |a||b|cos \alpha$
\item $[a, b] = x_1y_2 - x_2y_1= |a||b|sin \alpha$
\end{enumerate}

\subsection*{Уравнение прямой}
\addcontentsline{toc}{subsection}{Уравнение прямой}

В программировании мы будем задавать прямую с помощью двух точек. Воспользуемся уравнением ax+by+c=0. Плодить системы уравнений крайне сложно, поэтому мы будем стараться как можно реже сталкиваться с ними. Давайте возьмём какую-то точку (x, y) и составим уравнение, которое позволит нам сказать, что (x, y) лежит на этой прямой.

\begin{figure}[h]
\centering
\includegraphics[width=0.2\linewidth]{aaapic64.png}
\end{figure}

Какие вектора должны быть, чтобы (x, y) лежал на этой прямой? Коллинеарными. У коллинеарных векторов синус угла между ними равен 0. А значит, что вне зависимости от длины этих векторов, векторное произведение должно быть равно нулю. 

Давайте запишем x$_2$ - x$_1$ и y$_2$ - y$_1$. Это координаты вектора. Теперь применяем формулу: (x$_2$ - x$_1$)(y - y$_1$) - (x - x$_1$)(y$_2$ - y$_1$) = 0 и получаем x$_2$y - x$_2$y$_1$ - x$_1$y + x$_1$y$_1$ - xy$_2$ + xy$_1$ + x$_1$y$_2$ - x$_1$y$_1$ = 0. Выходит (y$_1$ - y$_2$)x + (x$_2$ - x$_1$)y + (-x$_2$y$_1$ + x$_1$y$_1$ + x$_1$y$_2$ - x$_1$y$_1$) = 0. Следовательно, чиселки перед х и у --- это a и b в уравнении соответственно. 

\subsection*{Вектор нормали}
\addcontentsline{toc}{subsection}{Вектор нормали}

Давайте найдём вектор нормали прямой, которой уже задана уравнением, приложенный к какой-нибудь точке. Итак, вам задана прямая ax+by+c=0. Вам нужно найти две какие-нибудь точки, лежащие на этой прямой. Идея <<подставь любой икс>> работает почти всегда, кроме вертикальной прямой. Можно \textit{рассмотреть случаи:} если b $\neq$ 0, то мы можем взять x = 0 или 1 и получить какие-то точки, а если b = 0, мы можем взять y = 0 или 1. 

Теперь, имея две точки, мы можем определить направляющий вектор. Обозначим его x$_p$, y$_p$, а вектор нормали как x, y.

\begin{figure}[h]
\centering
\includegraphics[width=0.3\linewidth]{aaapic65.png}
\end{figure}

Какой критерий того, что два вектора перпендикулярны друг к другу? x$_p$ $\cdot$ x + y$_p$ $\cdot$ y = 0, если вектор не нулевой. Сделать это можно, если записать x = -y$_p$ и y = x$_p$. Осталось найти направляющий вектор нашей прямой. Это -b,a (a,b), потому что a(x-b) + b(y+a) + c = 0, где ab сокращается. 

Как найти вектор нормали к нему? Мы берем x = -y$_p$, то есть -a, а y = x$_p$, то есть -b. Но вектор a, b тоже сгодится. 

\subsection*{Расстояние от точки до прямой}
\addcontentsline{toc}{subsection}{Расстояние от точки до прямой}

Есть прямая и точка. Необходимо посчитать расстояние от точки до прямой. 

\begin{figure}[h]
\centering
\includegraphics[width=0.4\linewidth]{aaapic66.png}
\end{figure}

Надо найти длину перпендикуляра, исходящего из точки xy. Построим треугольничек. 

\begin{figure}[h]
\centering
\includegraphics[width=0.4\linewidth]{aaapic67.png}
\end{figure}

Мы можем найти его площадь с помощью векторного произведения через боковые стороны и синус угла между ними. Потом делим полученную площадь на длину основания и получаем высоту треугольника, которая ещё и перпендикуляр. Теперь нам надо определиться, куда идти. Длину боковой стороны и высоты мы знаем, поэтому ищем длину противолежащего катета. А затем к направляющему вектору берем точку x$_1$y$_1$ и прикладываем заданной длины соответствующий отрезок.

\begin{figure}[h]
\centering
\includegraphics[width=0.4\linewidth]{aaapic68.png}
\end{figure}

Первая проблема заключалась в том, что мы могли пойти не в ту сторону. Мы её решили. Вторая проблема показана на картинке снизу.

\begin{figure}[h]
\centering
\includegraphics[width=0.2\linewidth]{aaapic69.png}
\end{figure}

В данном случае мы бы искали верно высоту. А также вектор нормали может смотреть вверх, а не в сторону прямой. Делая домашнее задание, пожалуйста, не забывайте, что векторы могут смотреть в разные стороны. 

\subsection*{В каком ухе жужжит?}
\addcontentsline{toc}{subsection}{В каком ухе жужжит?}

Фрекен Бок у нас представляется вектором. Слева у неё расположено левое ухо, а справа --- правое. И где-то в точке xy находится Карлсон. 

\begin{figure}[h]
\centering
\includegraphics[width=0.2\linewidth]{aaapic70.png}
\end{figure}

Как это решать? Обозначаем наши векторы и ищем знак синуса угла между ними. Если угол оказался меньше $\pi$, то Карлсон жужжит в левом ухе, а если от $\pi$ до $2\pi$, то в правом.  

\subsection*{Взаимное расположение фигур}
\addcontentsline{toc}{subsection}{Взаимное расположение фигур}

У нас есть две прямые: одна задана координатами x$_1$y$_1$ и x$_2$y$_2$, вторая --- x$_3$y$_3$ и x$_4$y$_4$. 

\begin{figure}[h]
\centering
\includegraphics[width=0.25\linewidth]{aaapic71.png}
\end{figure}

Нам нужно понять взаимное расположение прямых. Они могут пересекаться, совпадать и быть параллельными. Как проверить, что векторы коллинеарны? Их векторное произведение равно нулю. Чтобы разделить случаи того, что векторы параллельны и совпадают? Масса способов есть. Можно одну из точек одной прямой подставить в уравнение прямой другой, если у вас задано уравнениями. Если у вас задано точками, то можно подсчитать векторное произведение векторов из точки 1 в точку 2 и из точки 1 в точку 3, которое должно быть равным 0. Это означает, что точка 3 лежит на прямой x$_1$x$_2$, а еще мы знаем, что прямые коллинеарны. Следовательно, они совпадают.

Таким образом, мы разделили все три случая. Вернёмся к нашей задаче: найти точку пересечения двух прямых. Решать систему линейных уравнений мы не будем. Пронумеруем точки, чтобы не ссылаться на их координаты каждый раз:

\begin{figure}[h]
\centering
\includegraphics[width=0.27\linewidth]{aaapic72.png}
\end{figure}

Точкой пересечения будет точка Р. Строим подобные треугольники.

\begin{figure}[h]
\centering
\includegraphics[width=0.27\linewidth]{aaapic73.png}
\end{figure}

Теперь мы знаем не только высоты, но и их отношение. Зная высоты, мы можем сказать, что треугольники подобны по двум углам, и если у них высоты соотносятся каким-то образом между собой, то и все остальные стороны тоже. Следовательно, мы получаем пропорцию, где мы должны разделить наш отрезок между точками 3 и 4 и поставить нашу точку пересечения. Исходя из наших рассуждений, мы получили, что это просто отношение площадей треугольников 124 и 123. 

Но может возникнуть следующая проблема:
\newpage
\begin{figure}[h]
\centering
\includegraphics[width=0.27\linewidth]{aaapic74.png}
\end{figure}

Значит, площадь нужно считать ориентированную. Смотрите, в первом случае, считая косым произведением, 13 к 12 и 12 к 14 у нас получается два знака минус и они сокращаются. Во втором же случае 13 к 12 и 12 к 14 будет со знаком минус (см. картинку: точки 3 и 4 на одной стороне). Во втором случае отношение высот получится отрицательным, а значит, мы вектор от точки 4 к точке 3 домножаете на отрицательное число и он разворачивается как раз таки в ту сторону, где находится точка пересечения. 

В случае косого произведения вы получаете ориентированную площадь, а значит, модуль брать от неё не надо.

\subsection*{Многоугольники}
\addcontentsline{toc}{subsection}{Многоугольники}

Многоугольник задаётся путём обхода вершин. Как посчитать его площадь? 

\begin{figure}[h]
\centering
\includegraphics[width=0.24\linewidth]{aaapic75.png}
\end{figure}

Предположим, что у нас точка начала координат расположена где-то внизу. Пускай обход идёт по часовой стрелке.

\begin{figure}[h]
\centering
\includegraphics[width=0.24\linewidth]{aaapic76.png}
\end{figure}

Знак полученного треугольника --- это минус. Далее следующий треугольник со знаком плюс и так далее.

\begin{figure}[h]
\centering
\includegraphics[width=0.24\linewidth]{aaapic77.png}
\end{figure}

Вся площадь многоугольника учитывается. А кусочки со знаком минус сократились. Мы просто прошлись по всем точкам, составляющим границу многоугольника, и посчитали сумму косых произведений треугольников с вершиной в точке 0,0, потому что там удобнее считать, и всех отрезков, составляющих границу нашего многоугольника. Что самое удивительное, так это то, что это работает вне зависимости от того, где расположена точка 0,0. 

\subsection*{Биссектриса}
\addcontentsline{toc}{subsection}{Биссектриса}

Найдём угол, а потом поделим его на два. Как найти угол между двумя векторами? Ясно, что он должен быть $\leq \pi$. Представляю вам секретное знание: геометрическая функция atan2. Она принимает два параметра: первое --- некоторая константа С, умноженная на синус, второе ---- та же самая константа, умноженная на косинус. 

А как найти биссектрису-то? Сначала нужно нормализовать векторы, а потом просто сложить их.

\subsection*{Окружность. Поиски точек касания}
\addcontentsline{toc}{subsection}{Окружность. Поиски точек касания}

Окружность задаётся центром и радиусом. А также задана точка. Необходимо определить все точки касания, то есть уравнение касательной для нашей окружности. 
\begin{enumerate}
\item Как понять, что точка внутри окружности? Расстояние между точкой и центром окружности меньше радиуса. 
\item Как понять, что точка лежит на плоскости? Расстояние равно радиусу. 
\item Как понять, что точка вне плоскости? Расстояние больше радиуса.
\end{enumerate}

\begin{figure}[h]
\centering
\includegraphics[width=0.4\linewidth]{aaapic78.png}
\end{figure}

Если у нас есть точка касания, мы строим к ней перпендикуляр, соответственно мы знаем, насколько его нужно построить (в общем, всё про него знаем). Поэтому нам нужно отложить от этой точки вектор перпендикулярный заданному. Как это сделать? Нам нужно взять уравнение прямой, заданной двумя точками, и мы знаем, какой вектор нормали у неё --- а,b --- а потом от этой точки отложить вектор a,b и при необходимости посчитать уравнение. 

Давайте решать задачу в общем случае, когда у нас две касательные: найти точки касания. Перпендикуляр к касательной равен радиусу, поэтому мы можем посчитать гипотенузу (это расстояние между центром окружности и точкой). Значит, мы знаем длину касательной, радиус и расстояние между центром окружности и точкой. 

Конечно, можно поворачивать вектор с помощью матрицы, но я могу опустить перпендикуляр от точки касания до гипотенузы и получить некоторую точку В.

\begin{figure}[h]
\centering
\includegraphics[width=0.4\linewidth]{aaapic79.png}
\end{figure}

Через площади мы узнаем длину перпендикуляра. А также мы знаем, что треугольники ABC и ABP подобны. Раз уж я знаю всё про эти треугольники, значит, я могу посчитать длину отрезка ВР. Она равна отношению РА к РС. 

Получаем следующее равенство:

\[
|PB| = |PA| \cdot \frac{|PA|}{|PC|}
\]

А дальше откладываем вектор от точки Р в сторону точки С. Мы нашли точку В. А дальше мы считаем вектор нормали прямой РС, зная длину АВ (мы её знаем по теореме Пифагора через РВ и АР). Поэтому вектор нормали откладываем в одну сторону и не забываем отложить этот вектор нормали в другую сторону. И таким образом находим точки касания. 

\subsection*{Пересечение двух окружностей}
\addcontentsline{toc}{subsection}{Пересечение двух окружностей}

Точек пересечения может быть 2, 1, 0 и бесконечно много.
\begin{enumerate}
    \item Как понять, что их бесконечно много? Радиус и центр совпадают.
    \item Когда одна точка пересечения? Когда сумма радиусов равна расстоянию между центрами окружностей и внутреннее касание (когда расстояние между центрами + радиус внутренней окружности = радиусу другой окружности).
    \item Когда у них 0 точек касания? Когда расстояние между центрами + радиус внутренней окружности < радиуса другой окружности, а также сумма радиусов меньше расстояния между центрами окружностей.
    \item Когда у нас две точки касания? Нужно учесть внутренние и внешние касания одновременно, а для этого нужно перестать брать модуль без надобностью.
\end{enumerate} 

Как только вы берёте модуль, у вас появляются случаи. Давайте искать:

Мы знаем все три стороны в нашем прекрасном треугольнике (радиусы и расстояния между центрами).

\begin{figure}[h]
\centering
\includegraphics[width=0.4\linewidth]{aaapic80.png}
\end{figure}

Что мы можем еще посчитать? Высоту, зная основание и площадь. Нам осталось найти точечку А.

\begin{figure}[h]
\centering
\includegraphics[width=0.4\linewidth]{aaapic81.png}
\end{figure}

Теперь мы знаем все три стороны и высоту. Мы можем довольно легко посчитать длину отрезка от центра левой окружности до точки А. А дальше откладываем вектор от центра левой окружности к центру правой окружности. 

\subsection*{Как это всё писать?}
\addcontentsline{toc}{subsection}{Как это всё писать-то?}

Я вам советую написать как минимум функции. НО если вы хотите совсем хорошо это сделать, то я вам советую писать классы с методами (в геометрии это очень пригождается). А также ни в коем случае не решайте систему уравнений. 
\newpage

\section*{Лекция 10: Продвинутая геометрия}
\addcontentsline{toc}{section}{Лекция 10: Продвинутая геометрия} 

\subsection*{Динамика по профилю}
\addcontentsline{toc}{subsection}{Динамика по профилю}

Идея: у нас была динамика по подмножествам, когда мы решали с вами задачки про коммивояжёра. Давайте я вам напомню. Простое решение имеет сложность порядка N!, потому что мы перебираем всевозможные пути. Как мы это оптимизировали: у нас есть динамика, у которой есть параметр (подмножество) городов, которые коммивояжёр уже посетил, и в каком городе он сейчас находится. Как эту динамику посчитать через меньшую динамику? Как понять, что он уже был в subset'е городов и сейчас находится в городе last? Поэтому мы просто переберём все варианты предпоследнего города, где он был. То есть это:

\[
d[subset, last] = \min_{prev \in subset/(last)}  d[subset/(last), prev] + g[prev, last]
\]

Мы перебираем подмножество городов, где у нас коммивояжёр находится, в каком из городов этого подмножества он находится сейчас. Как мы можем попасть в это состояние? Мы оптимальным образом посетили всё подмножество, за исключением этого самого последнего города. Перебираем всевозможные предпоследние города, которые принадлежат собственно вот этому подмножеству, исключая последний, и прибавляем длину ребра между предпоследним и последним городом. Таким образом, насчитываем. 

Тут нужно подумать, как нам хранить подмножество. Хранить мы его будем с помощью битовой маски. Сложность этого алгоритма будет O(2$^N$ $\times$ N$^2$). 

Вернёмся к динамике по профилю. 

Решим задачку: количество способов замастить некоторый прямоугольник N $\times$ M доминошками 1 на 2. Чтобы жить мёдом не казалась, есть места на доске, где доминошку ставить нельзя. Поле будет не большое.

В терминах двудольного графа нам надо посчитать количество паросочетаний. Но мы не будем к этому сводить и попытаемся решить так. Сделаем динамику следующего рода: пусть мы считаем k-ый столбец. Всё, что левее его, уже полностью корректно заполнено доминошками. Что может произойти? У нас из k-1-го столбца в k-ый могут торчать какие-то доминошки, положенные горизонтально.

\begin{figure}[h]
\centering
\includegraphics[width=0.4\linewidth]{aaapic82.png}
\caption{Зачёркнутые квадраты --- выколотые, а поперёк стоящие прямоугольники --- доминошки}
\end{figure}

Идея, как это считать: k-ый столбец мы назовём профилем, по которому идёт динамика. И нас может быть много вариаций, как этот столбец выглядит. Например, мы какие-то доминошки можем поставить вертикально. То есть для каждого профиля будет своё количество вариантов, сколько можно получить в k-ом столбце такой профиль. Профиль --- это состояние последнего столбца. Кодируется оно точно так же ноликами и единичками.

Теперь нам нужно научиться переходим к следующему столбцу, от нашего состояния вперед. То есть посчитать наше состояние через предыдущее. Здесь нам снова понадобится динамика: номер столбца, для которого мы считаем, и профиль.

\[
dp[k, profile]
\]

Условно, пусть у нас профиль предыдущего столбца был такой:\\
0 \\
1 \\
0 \\
1 \\
1 \\

К каким профилям мы можем перейти, чтобы соблюдалось наше правило, что все предыдущие столбцы заполнены? Это означает вот такую ситуацию:

\newpage
\begin{figure}[h]
\centering
\includegraphics[width=0.2\linewidth]{aaapic83.png}
\caption{Это так доминошки торчат}
\end{figure}

К каким профилям мы можем перейти?

\begin{figure}[h]
\centering
\includegraphics[width=0.2\linewidth]{aaapic84.png}
\end{figure}

Нам нолик обязательно нужно покрыть, сделать это можно указанным на картинке способом (в коде нужно писать двумя циклами). Конечно, зависеть это будет от того, как вы будете писать. То есть мы в предыдущий столбец не кладём доминошки, чтобы не считать одно и то же несколько раз. А если у нас несколько нулей подряд, то далее мы можем расположить вертикальную доминошку.

\begin{figure}[h]
\centering
\includegraphics[width=0.25\linewidth]{aaapic85.png}
\end{figure}

Соответственно, что нам нужно сделать? Нам нужно сделать матричку совместимых профилей. Что мы можем от профиля 1 до профиля 2 корректно перейти, не оставляя никаких дырок и не накладывая доминошку дважды. 
\newpage
\begin{figure}[h]
\centering
\includegraphics[width=0.4\linewidth]{aaapic86.png}
\end{figure}

от 26-го профиля (2+8+16) мы можем перейти к 5-ому (1+4), значит, в пятом профиле ставим единичку. А ещё к 29-ому (1+4+8+16).

\begin{figure}[h]
\centering
\includegraphics[width=0.4\linewidth]{aaapic87.png}
\end{figure}

В остальных клеточках у нас нули. Это вариант пока без выколотых точек. Как эту табличку построить? Перебрать всевозможные варианты. Проблема в чём? С горизонтальной понятно: у нас все нули должны превратиться в единицы в совместимом профиле. А вот там, где у нас было несколько подряд единиц в предыдущем профиле, мы должны перебрать всевозможные варианты расставить доминошки размера 2. Ограничения на N в реальности не больше 10, поэтому это можно сделать практически как угодно. 

Как посчитать динамику в заданном профиле? Нужно просто просуммировать количество вариантов.

\[
dp[k, profile] = \sum_{pr2 = 0}^{2^{N-1}} dp[k-1, pr2] \times m[pr2, profile]
\]

А если мы добавляем выколотые клетки? Мы как бы заполняем их. То есть в этом профиле обязаны стоять единички. А в предыдущем профиле тоже обязана стоять единичка, потому что мы договорились, что вертикальные доминошки мы размещаем только в текущем столбце. А все остальные профили мы игнорируем. 

\subsection*{Симпатичные узоры}
\addcontentsline{toc}{subsection}{Симпатичные узоры}

У нас есть табличка N $\times$ M $\leq$ 30. Надо расположить плитки так, чтобы узор был симпатичный. Симпатичным называется узор, в котором нет ни одного квадратика 2 на 2 полностью белого или полностью чёрного цвета. Нам нужно посчитать количество симпатичных узоров для заданного размера N $\times$ M. 

Идея абсолютно такая же: посчитаем динамику по профилю. Всё, что слева k-го профиля, мы раскрасили правильно. Единички и нолики мы должны располагать так, чтобы не возникало несимпатичного узора. 

Давайте напишем все профили высотой 2:

\begin{figure}[h]
\centering
\includegraphics[width=0.3\linewidth]{aaapic88.png}
\end{figure}

Как выглядит его матрица совместимости? Здесь у нас всё совместимо, кроме таких маленьких профилей, кроме полностью белых или полностью чёрных квадратиков.

Что мы делаем? Когда мы считаем динамику для очередного столбца и очередного профиля, мы просто перебираем предыдущий столбец (k-ый), всевозможные профили и домножаем на матрицу. 

Какой прикол в этой задаче? Поскольку у нас стоит ограничение N $\times$ M $\leq$ 30, нам нужно было правильно ориентировать стороны нашего многоугольника. Понятно, что количество способов раскрасить прямоугольники 5 на 6 и 6 на 5 --- одинаково. Наш наихудший случай, условно, это попытка сделать 15 на 2, потому что мы просто не уложимся. Поэтому за количество строк мы берем наименьшее из чисел. 

За какую сложность у нас всё это работает? Чтобы посчитать динамику для столбца и профиля, нам нужно перебрать все предыдущие профили. Получаем M $\times$ 2$^N$ $\times$ 2$^N$. 

Теперь пускай у нас N $\leq$ 5 и M $\leq$ 10000000. Давайте придумаем, как нам организовать переход от предыдущего профиля через всевозможные профили к нашему профилю настоящему. 

\begin{figure}[h]
\centering
\includegraphics[width=0.4\linewidth]{aaapic89.png}
\end{figure}

То есть мы хотим посчитать не просто совместимость двух соседних профилей, а совместимость двух профилей через всевозможные промежуточные. Понятно, что кол-во вариантов перейти от одного профиля через всевозможные промежуточные к какому-то одному профилю может быть больше одного. Потому что может быть несколько промежуточных профилей, которые совместимы с предыдущим и нынешним. Как мы это можем посчитать? Мы можем посмотреть на всевозможные средние профили, и если этот средний профиль одновременно совместим и с предыдущим, и со следующим, то добавляем единичку к количеству таких переходов. И можем составить матрицу переходов через один шаг. 

На самом деле все эти вышеуказанные профили --- это матрица, возведённая в квадрат. Когда мы возводим матрицу в квадрат, ij-ая ячейка будет считаться как произведение i-ой строки на j-ый столбец. А какой физический смысл в этом? Что мы, возведя эту матрицу в квадрат, получим матрицу перехода через один промежуточный профиль. Почему так получается? В нашей новой матрице, допустим, мы хотим посчитать количество переходов из 00 в 01:

\begin{figure}[h]
\centering
\includegraphics[width=0.3\linewidth]{aaapic90.png}
\end{figure}

Выглядит это так: m2[00, 01]. Что мы делаем? Мы перебираем первый промежуточный столбец, всевозможные промежуточные столбцы и смотрим произведение первого элемента в строке на первый элемент в столбце: 1 $\cdot$ 0 = 0. Какой физический смысл у происходящего? Мы попытались попасть из профиля 00 в профиль 01 через промежуточный профиль 00 промежуточный (смотрите, если непонятно: из 00 в 01 --- это то, что сбоку, а промежуточный 00 --- это тот, что сверху). Из 00 в 01 подать можно, а из 00 в 00 нельзя, потому что существует 0 вариантов. Теперь мы пытаемся попасть через промежуточный столбец 01: у нас 1 способ попасть из 00 в 01 и 1 способ попасть из 01 в 01. Перемножаем, получаем 1. Теперь пытаемся попасть через 10. Из 00 в 10 можем попасть, из 10 в 01 тоже можем попасть. Посчитали произведение, единичка. Теперь через промежуточный столбец 11: из 00 в 11 попадать можно, из 11 в 01 попадать можно. Итого оказалось 3 способа попасть из 00 в 01 через всевозможные промежуточные столбцы.

А как ходить через два промежуточных столбца? Чтобы ходить через один промежуточный столбец, нужно возвести матрицу в квадрат, а чтобы через два, то нужно возвести её в куб. Почему это работает? Потому что если мы эту квадратную матрицу умножим ещё раз на исходную матрицу, то получится количество способов перейти через два промежуточных (в нашем случае мы ещё один вариант добавим). 

Получается, чтобы построить матрицу перехода на M шагов, нам нужно исходную матрицу возвести в M-1-ую степень. Объекты мы умеем возводить в степень за логарифм с помощью быстрого возведения в степень. Размер матрицы у нас 2$^{2N}$, а также произведение у нас не мгновенное. Таким образом, мы получаем сложность 2$^{3N}$logM. 

\subsection*{Цветная лента}
\addcontentsline{toc}{subsection}{Цветная лента}

Профилем не всегда бывает столбец. Есть длинная полоска, высотой 1 и длиной N клеток. Каждая плиточка покрашена в красный, синий и зеленый цвет.

\begin{center}
\begin{tabular}{|c|c|c|c|c|}
\hline
к & c & з & з & \ldots \\
\hline
\end{tabular}
\end{center}

Требуется определить кол-во вариантов раскрасить полоску длиной N в красные, синие, зеленые цвета так, чтобы у неё была заданная матрица соседних цветов. Что это? Если у нас есть пара соседних клеток, например, красный и синий, то мы помечаем их в таблице единичкой. Вот картиночка:

\begin{center}
\begin{tabular}{|c|c|c|c|}
\hline
- & к & с & з \\
\hline
к & 0 & 1 & 0 \\
\hline
с & 1 & 0 & 1  \\
\hline
з & 0 & 1 & 1 \\
\hline
\end{tabular}
\end{center}

Получаем матрицу соседей. Нужно определить количество полосок длины N, с заданной матрицей соседей. 

Что у нас будет являться профилем в данном случае? В динамике один из ее параметров будет длина полоски, а второй --- профиль. Мы что-то знаем про полоски длиной k$-$1, и сейчас мы хотим поставить какую-нибудь плиточку, чтобы всё было хорошо и узнать ответ для полосок длиной k. Как проверить, что всё хорошо? Допустим, мы добавляем красную клетку. Что произойдёт? Всё сломается, если на место нуля появится единичка. По нашей вышеуказанной матрице два подряд красных были запрещены, поэтому, если в полоске предпоследняя была красная, то мы не можем приделать к ней еще одну красную плитку, потому что это будет противоречить нашей матричке. Поэтому к динамике мы сейчас добавим последний параметр: last color. А чтобы профиль не смущал нас, я переименую его в матрицу. На самом деле матрица целиком мне не нужна. Сколько элементов матрицы мне нужно? Я могу выбросить все элементы ниже главной диагонали, потому что матрица симметрична. То есть, на самом деле, мне нужно только 6 битиков тоскать в нашей матрице, ведь они ее и задают. Я хочу посчитать количество полосок, длинной k, заданной матрицей и с известным последним цветом плитки. Как мне это посчитать? Откуда я могу попасть в это состояние? Понятно, что я могу попасть только из длины k$-$1. То есть это:

$$
\begin{align}
dp[k, matrix, last\_color] = \sum_{prev\_color \in \text{(к, с, з)}}^{(color, prev\_color) \in matrix} dp[k-1, matrix, prev\_color] \\
+ \sum_{prev\_color \in \text{(к, с, з)}}^{\substack{(color, prev\_color) \in matrix \\ matr' = matrix/(color, prev\_color)}} dp[k-1, matr', prev\_color]
\end{align}
$$

Мы записываем сумму, где цвет и его предыдущий принадлежит матрице, но не принадлежит предыдущей матрице matr', где matr' = matrix/(color, prev\_color). Это значит, у нас пара (color, prev\_color) встретилась впервые, поэтому нам нужна матрица со штрихом, у которой на месте нуля стояла единичка. А также обязательно проверять, что пара (color, prev\_color) содержится в матрице.

\subsection*{Литограф для плат}
\addcontentsline{toc}{subsection}{Литограф для плат}

Инженеры решили начать делать платы, но некоторые блоки оказались битые. Нам нужно навырезать как можно больше плиточек размером 2 $\times$ 3. 

Здесь нам понадобится троичный профиль (о боже). Его идея: у нас в каждый момент времени будет происходить следующее. Мы про профиль знаем, что уже в столбце занята первая клетка. Это был первый вариант. Второй вариант --- это когда вторая клетка столбца пуста и слева от неё заполнена 1 клеточка. Третий вариант --- это когда у нас две клеточки свободны: в нынешнем и предыдущем столбце (см. картинку).

\begin{figure}[h]
\centering
\includegraphics[width=0.2\linewidth]{aaapic91.png}
\caption{Профиль задаётся чиселками 0, 1 и 2}
\end{figure}

Когда мы можем разместить блок горизонтально? Когда три предыдущие единички или двоечки. А чтобы горизонтально вырезать чип? Из какого состояния мы можем перейти? Из двух двоек мы переходим в состояние двух нулей, вырезая чип. 

Таким образом, мы рассмотрели основные алгоритмы и структуры данных, необходимые для эффективного решения задач. Но алгоритмы, не вошедшие в этот курс и лишь бегло упомянутые в конце последней лекции, – это уже совсем другая история.






\newpage
всё
\newpage
передаю привет своим самым преданным читателям и...
\newpage
благодарю
\newpage
\[
\text{За всё самое хорошее}
\]
\begin{figure}[ht]
\centering
\begin{minipage}{0.45\textwidth}
  \centering
  \includegraphics[width=\textwidth]{image1.jpg}
  \label{fig:img1}
\end{minipage}
\hfill
\begin{minipage}{0.45\textwidth}
  \centering
  \includegraphics[width=\textwidth]{image2.jpg}
  \label{fig:img2}
\end{minipage}
\end{figure}
\newpage
\[
\text{И всё самое трудное}
\]
\begin{figure}[h]
\centering
\includegraphics[width=0.5\linewidth]{image3.jpg}
\end{figure}
\newpage
Неплохая получилась история... интересная, веселая, порой немного грустная, а самое главное --- поучительная. Она научила нас быть смелыми и не бояться вызовов, которые готовит нам жизнь, помогала добиваться нам поставленных целей, несмотря ни на что. А что самое важное, она научила нас по-настоящему любить и не сходить с пути, следуя за своей мечтой. Встретимся на экзамене.
\newpage
\\
\\
\\
\\
\\
\\
\begin{center}
\fontsize{70}{80}\selectfont\textbf{КОНЕЦ}
\end{center}
\begin{center}
\fontsize{15}{20}\selectfont{А кто слушал --- молодец)}
\end{center}





\end{document}